\documentclass[11pt]{article}

\usepackage[margin=1.02in]{geometry}
\usepackage{amsmath,amssymb,amsthm,mathtools,mathrsfs}
\usepackage{bm}
\usepackage{enumitem}
\usepackage{booktabs,array,tabularx}
\usepackage{microtype}
\usepackage[colorlinks=true,linkcolor=blue,citecolor=blue,urlcolor=blue]{hyperref}
\usepackage[nameinlink,capitalise]{cleveref}

% -----------------------------------------------------------------------------
% separate tables of contents for the main text and Mathematical Supplement
% -----------------------------------------------------------------------------
\newif\ifmathsupplement
\mathsupplementfalse

\makeatletter
\let\RG@addcontentsline\addcontentsline
\renewcommand{\addcontentsline}[3]{%
  \def\RG@tocname{#1}%
  \def\RG@toc{toc}%
  \ifx\RG@tocname\RG@toc
    \ifmathsupplement
      \RG@addcontentsline{supptoc}{#2}{#3}%
    \else
      \RG@addcontentsline{maintoc}{#2}{#3}%
    \fi
  \else
    \RG@addcontentsline{#1}{#2}{#3}%
  \fi
}

\newcommand{\maintableofcontents}{%
  \section*{Contents}%
  \IfFileExists{\jobname.maintoc}{%
    \@starttoc{maintoc}%
  }{%
    \newwrite\RGmaintoc
    \immediate\openout\RGmaintoc=\jobname.maintoc
    \immediate\closeout\RGmaintoc
    \@starttoc{maintoc}%
  }%
}

\newcommand{\supplementtableofcontents}{%
  \section*{Contents of the Mathematical Supplement}%
  \IfFileExists{\jobname.supptoc}{%
    \@starttoc{supptoc}%
  }{%
    \newwrite\RGsupptoc
    \immediate\openout\RGsupptoc=\jobname.supptoc
    \immediate\closeout\RGsupptoc
    \@starttoc{supptoc}%
  }%
}

\makeatother
\usepackage{tikz-cd}
\setlength{\emergencystretch}{4em}
\raggedbottom
\hbadness=10000

% -----------------------------------------------------------------------------
% theorem environments
% -----------------------------------------------------------------------------
\newtheorem{mastertheorem}{Master Theorem}
\newtheorem{theorem}{Theorem}[section]
\newtheorem{proposition}[theorem]{Proposition}
\newtheorem{lemma}[theorem]{Lemma}
\newtheorem{corollary}[theorem]{Corollary}
\newtheorem{countertheorem}[theorem]{Countertheorem}
\newtheorem{example}[theorem]{Example}
\theoremstyle{definition}
\newtheorem{definition}[theorem]{Definition}
\newtheorem{construction}[theorem]{Construction}
\newtheorem{assumption}[theorem]{Assumption}
\theoremstyle{remark}
\newtheorem{remark}[theorem]{Remark}
\newtheorem{interpretation}[theorem]{Interpretation}
\newtheorem{scope}[theorem]{Scope statement}

\crefname{mastertheorem}{Master Theorem}{Master Theorems}
\Crefname{mastertheorem}{Master Theorem}{Master Theorems}

% -----------------------------------------------------------------------------
% notation
% -----------------------------------------------------------------------------
\newcommand{\N}{\mathbb N}
\newcommand{\Z}{\mathbb Z}
\newcommand{\R}{\mathbb R}
\newcommand{\C}{\mathbb C}
\newcommand{\Torus}{\mathbb T}
\newcommand{\one}{\mathbf 1}
\newcommand{\id}{\mathrm{id}}
\newcommand{\dd}{\mathrm d}
\newcommand{\Tr}{\operatorname{Tr}}
\newcommand{\tr}{\operatorname{tr}}
\newcommand{\rank}{\operatorname{rank}}
\newcommand{\Ker}{\operatorname{Ker}}
\newcommand{\Ran}{\operatorname{Ran}}
\newcommand{\Span}{\operatorname{span}}
\newcommand{\Hom}{\operatorname{Hom}}
\newcommand{\End}{\operatorname{End}}
\newcommand{\supp}{\operatorname{supp}}
\newcommand{\spec}{\operatorname{spec}}
\newcommand{\diag}{\operatorname{diag}}
\newcommand{\norm}[1]{\lVert#1\rVert}
\newcommand{\abs}[1]{\lvert#1\rvert}
\newcommand{\ip}[2]{\left\langle #1,#2\right\rangle}
\newcommand{\A}{\mathcal A}
\newcommand{\Hh}{\mathcal H}
\newcommand{\Om}{\Omega}
\newcommand{\GT}{\boldsymbol{\mathbb T}_{\mathrm{Grand}}}
\newcommand{\GTX}{\boldsymbol{\mathbb T}_{\mathrm{Grand},X}}
\newcommand{\Lip}{\operatorname{Lip}}
\newcommand{\Cl}{\mathrm{Cl}}
\newcommand{\HS}{\mathrm{HS}}
\newcommand{\vol}{\operatorname{vol}}
\newcommand{\supplement}{Appendices}
\newcommand{\EssIm}{\operatorname{EssIm}}
\newcommand{\Curr}{\operatorname{Curr}}
\newcommand{\RReal}{\operatorname{RReal}}
\newcommand{\Corr}{\operatorname{Corr}}
\newcommand{\Stab}{\operatorname{Stab}}
\newcommand{\McM}{\operatorname{McM}}
\newcommand{\RG}{\operatorname{RG}}
\newcolumntype{L}[1]{>{\raggedright\arraybackslash}p{#1}}


\title{\textbf{From Predictive Dynamics to Spectral Geometry}}

\author{%
  Aur\'elien P\'elissier$^{\dagger}$\\
  {\small $^{\dagger}$Correspondence: \href{mailto:aurelien.pelissier.38@gmail.com}{aurelien.pelissier.38@gmail.com}}%
  \vspace{-0.5cm}
}
\date{}



\begin{document}
\maketitle



\begin{abstract}
In Alain Connes' noncommutative geometry, geometry is encoded by an algebra, a Hilbert space, and a Dirac operator taken as fundamental spectral data. We show that these structures can instead emerge from a more primitive predictive dynamics. A compatible first-order variation on the induced history algebra canonically generates a real even Hodge--Dirac spectral triple, while the algebra alone determines neither its metric scale nor the full Dirac operator.
%
The resulting spectral geometry is an information-reduced description: dynamically inequivalent predictive systems can share the same effective spectral data. Nevertheless, finite geometries become reconstructible after finitely many calibrations, and on periodic three-dimensional lattice approximants the emergent spectral distance converges to a Riemannian metric reconstructed independently from the dynamics.
%
Thus Connes-type spectral geometry may arise not as fundamental structure, but as an effective geometric description of richer underlying dynamics. This provides a concrete setting in which noncommutative geometry need not be fundamental, but may instead arise as a coarse-grained spectral description.
%
A \texttt{Lean} formalization is available at
\url{https://github.com/AI-MathPhys/renewal-geometry}.

\end{abstract}

\medskip
\noindent\textbf{Keywords.}
Noncommutative geometry; spectral triples; Renewal Geometry; predictive dynamics;
Hodge--Dirac operators; Connes distance; emergent geometry; spectral reconstruction.
\setcounter{tocdepth}{2}


\vfill

\maintableofcontents


\newpage



% =============================================================================
\section{Introduction}
\label{sec:introduction}
% =============================================================================

Many spaces that arise naturally in mathematics and physics do not admit a
satisfactory description as ordinary geometric spaces. Quotients may be singular,
points may fail to capture the relevant observables, and the structures of interest
may be encoded more faithfully by operators than by a classical manifold. In such
situations the usual geometric language---built from points, coordinates, and
functions on a well-behaved space---can cease to be adequate, even though an
algebra of observables remains meaningful.

Noncommutative geometry was developed to address this problem. Connes' programme
replaces the requirement of an underlying point-set space by an operator-algebraic
description in which an algebra plays the role of generalized coordinates. In the
commutative case this remains compatible with ordinary geometry, since a suitable
commutative algebra determines the underlying space through Gelfand duality.
Allowing the algebra to become noncommutative extends the same principle to
settings for which no classical space of points need exist.

The spectral formulation equips this generalized coordinate algebra with geometry.
A compact spin manifold can be encoded by an algebra $\mathcal A$ represented on
a Hilbert space $\mathcal H$, together with a Dirac operator $D$
\cite{Connes1994}. The representation supplies the operator-theoretic carrier,
while $D$ encodes differential and metric information.  Under the reconstruction
axioms for commutative spectral triples, this operator data recovers compact spin
geometry \cite{RennieVarilly2008,connesreconstruction}; in the genuinely
noncommutative case the same operator-geometric language continues to apply even
when no ordinary underlying space exists.

Spectral triples consequently provide a common language for noncommutative spaces,
differential calculi, index theory, and metric geometry. The Connes distance
formula turns the commutator with $D$ into a metric notion, making the Dirac
operator simultaneously a differential and geometric object. Related work on
compact quantum metric spaces formulates metric information at the level of an
algebra equipped with a Lipschitz seminorm \cite{Rieffel1999}.  Rieffel's quantum
Gromov--Hausdorff distance and the later propinquity and spectral propinquity
provide frameworks for convergence of quantum metric spaces and metric spectral
triples \cite{Rieffel2004,Latremoliere2022,FarsiLatremolierePacker2024}, while real
structures and gradings encode analogues of spin and orientation data.

A different structural question is whether the objects entering this spectral
description must themselves be regarded as primitive. This is logically prior to
spectral reconstruction: rather than asking which geometry is encoded by a given
spectral triple, one asks whether the algebra, representation, differential
structure, and Dirac data can arise from an underlying dynamics in which none of
them is initially prescribed.

\paragraph{From predictive dynamics to spectral geometry.}

Renewal Geometry provides one setting in which this question can be posed
operationally \cite{PelissierRenewalSpacetime2026}. Its starting point is not a
manifold, graph, causal order, algebra, Hilbert space, or spectral triple, but a
finite operational process specified by composable intervention histories and
their conditional future statistics. Histories that cannot be distinguished by
any admissible future continuation are identified. The resulting
future-equivalence quotient reconstructs the minimum persistent predictive
carrier of the process rather than requiring such a carrier to be supplied
independently.

Complete operational statistics then reconstruct a canonical positive history
object, the \emph{Renewal Grand Tensor}.  We denote the Grand Tensor associated
with a finite predictive system $X$ by $\GTX$.  Its role is pregeometric: it
records which distinctions survive prediction, how represented histories compose,
and how they overlap, before positions, distances, metrics, or continuum geometry
have been assigned.  The corresponding history algebra and positive moment
structure therefore arise from the operational process rather than being inserted
as geometric input.

For the present paper, the relevant interface is the algebraic content carried by
$\GTX$.  The Grand Tensor canonically supplies the future-minimal predictive
carrier, the represented history algebra $\A_X^{\mathrm{hist}}$, and its positive
moment structure.  On the finite algebra there is in particular an intrinsic
normalized regular trace, which will provide one canonical choice of $\tau_X$.
To pass from this reconstructed algebraic object to spectral geometry, one then
specifies a compatible first-order variation $\delta_X$ of the history algebra.
Thus the interface used here is
\begin{equation}
\GTX
\longmapsto
(\A_X^{\mathrm{hist}},\tau_X)
\overset{\delta_X}{\longrightarrow}
(\A_X,\Hh_X,D_X,J_X,\gamma_X).
\label{eq:RGT-spectral-interface}
\end{equation}
The first arrow is reconstructed directly from the operational statistics and is
recalled self-containedly in Section~\ref{sec:finite-main}.  The second is the
Hodge--Dirac spectralization studied in this paper.  The Grand Tensor is not
identified with the displayed algebraic or spectral data: it retains the
predictive carrier and complete positive history structure, including information
that need not survive in the spectral image.  Nor do we assume that the Grand
Tensor or its algebra alone fixes $\delta_X$; the first-order variation is exactly
the additional datum whose metric scale and Dirac ambiguity are analyzed below.

The same $\GTX$ also supports the separate Lorentzian spacetime construction
developed within the Renewal Geometry framework~\cite{PelissierRenewalSpacetime2026}.
The present work asks instead whether
Connes-type spectral geometry can arise from a compatible first-order variation
of its reconstructed history algebra.  The spacetime-specific assumptions of the
companion construction are not used below; once a finite-dimensional unital
$C^*$-algebra $\A$, a faithful tracial state $\tau$, and a compatible real
Hilbert-bimodule derivation $\delta$ are supplied, the finite spectralization
results are ordinary operator-theoretic statements independent of the operational
route by which $\A$ was obtained.

The resulting problem has both a forward and an inverse aspect. On the forward
side one asks whether a first-order predictive differential canonically produces
the structures of spectral geometry, and whether the resulting metric has an
independent dynamical meaning. On the inverse side one asks which spectral
geometries can be reached, what additional data are required to reconstruct a
full Dirac operator, and whether different microscopic dynamics can have the same
spectral image.

\paragraph{What lies before a spectral triple?}

At the finite-dimensional level, Krajewski classified finite spectral triples in
terms of finite matrix data and diagrams \cite{Krajewski1998}.  This gives a
powerful classification of already specified finite spectral geometries, but it
does not address whether the algebraic carrier or Dirac data can themselves be
reconstructed from an operational dynamics.

Several constructions derive spectral triples from other supplied geometric or
hierarchical structure.  For approximately finite-dimensional $C^*$-algebras, a
Dirac operator can be built naturally from a prescribed finite-dimensional
filtration \cite{ChristensenIvan2006}.  Weighted trees, ultrametric Cantor sets,
fractal spaces, stationary Bratteli diagrams, and aperiodic systems provide further
constructions in which the underlying space, algebra, hierarchy, or incidence
structure is already given
\cite{PearsonBellissard2009,ChristensenIvanSchrohe2012,KellendonkSavinien2012}.
These results show that the Dirac operator need not be specified independently,
but their mathematical entrance remains a pre-existing algebraic or geometric
carrier.

A related dynamical literature constructs or extends spectral triples from a
prescribed $C^*$-dynamical system.  Crossed-product constructions attach spectral
data to given group actions \cite{Paterson2014}; for symbolic dynamics, spectral
triples have been constructed on subshift algebras and their crossed products
\cite{JulienPutnam2016}; and analogous constructions have been developed for
irreversible $C^*$-dynamical systems \cite{AielloGuidoIsola2022}.  Here again the
dynamics acts on an algebra that is supplied before spectralization.  The present
question differs at the entrance: the history algebra itself is reconstructed only
after predictive future-equivalence has removed operationally invisible
refinements.

A particularly close line of work comes from quantum Riemannian geometry. Majid
has shown that a geometrically realized Dirac operator can be constructed
canonically from a prescribed noncommutative coordinate algebra, differential
calculus, quantum metric, and compatible bimodule connection, with explicit
realizations on graphs, lattices, the fuzzy sphere, and matrix algebras
\cite{Majid2025}. This places the Dirac operator downstream of quantum metric and
differential data. The present question begins one level earlier: neither the
coordinate algebra nor the differential and metric structures from which the
Dirac operator is built are assumed at the operational entrance.

An even closer dynamical precedent is provided by noncommutative potential theory.
Cipriani and Sauvageot showed that a tracially symmetric Dirichlet form on a
$C^*$-algebra admits a canonical realization through a closable derivation into a
Hilbert bimodule, with the Markov generator represented by the corresponding
divergence \cite{CiprianiSauvageot2003}.  This calculus has been used to define
Dirac operators on fractals \cite{HinzTeplyaev2013}, and Hinz, Kelleher and
Teplyaev constructed spectral triples from regular Dirichlet and resistance forms,
showing in the compact strongly local case that the Connes distance coincides with
the intrinsic metric \cite{HinzKelleherTeplyaev2015}.  Subsequent work has related
Hodge--Dirac spectral triples of sub-Markovian semigroups to heat-kernel and
spectral-dimension estimates \cite{Cipriani2016,Arhancet2024}.  These constructions
are especially close to the Hodge--Dirac layer used below, but they begin with the
algebra, reference state or trace, and Dirichlet or Markov structure already
supplied.

The operational entrance also has established precedents. Quantum combs and
process-tensor formalisms encode complete multitime intervention statistics by
positive causal process objects and describe memory-bearing quantum dynamics
without imposing a Markov approximation
\cite{ChiribellaDArianoPerinotti2009,PollockProcessTensor2018,PollockEtAl2018}.
They reconstruct or represent a multitime process from intervention statistics,
but do not by themselves identify histories by predictive future-equivalence and
then use that quotient to construct the coordinate algebra of a spectral geometry.

There is also a broad classical literature in which state is defined through
future predictions rather than hidden microscopic labels.  Observable operator
models encode stochastic processes through operators acting on observable
statistics \cite{Jaeger2000}, while predictive-state representations use
multi-step, action-conditional predictions as state coordinates
\cite{LittmanSuttonSingh2001,SinghJamesRudary2004}.  In computational mechanics,
causal states identify pasts that induce the same conditional distribution over
futures and yield minimal predictive presentations
\cite{CrutchfieldYoung1989,shalizicrutchfield,TraversCrutchfield2025}; continuous-
time extensions include renewal and discrete-event processes
\cite{MarzenCrutchfield2017}.  Renewal Geometry combines this predictive quotient
with multitime operational composition and positive process structure.  The
persistent carrier and represented history algebra are reconstructed only after
future-invisible distinctions have been removed, and the first-order differential
is then spectralized on that reconstructed algebra.

The question considered here can therefore be stated as follows:
\begin{quote}
\emph{Can the algebra, Hilbert-space carrier, differential structure, and Dirac
data of noncommutative geometry arise from an underlying predictive dynamics in
which none of them is initially prescribed?}
\end{quote}

A nontrivial answer requires more than reproducing a desired spectral triple by
definition. The algebra cannot simply be hidden in the microscopic model under
another name, and the Dirac operator cannot be fitted afterward to recover a
target metric. One must determine what the differential actually fixes, which
spectral geometries are obtainable under different notions of reconstruction,
whether the resulting metric can be validated independently, and what information
is lost in the passage from dynamics to spectral geometry.

The last point is central. If distinct predictive systems have the same spectral
image, then a spectral triple is not, in general, a complete invariant of the
dynamics from which it arose. It may instead describe an equivalence class of
microscopically different systems that agree on their effective metric and
differential structure. In this sense spectral geometry can be
\emph{information-reduced}: complete for the operator-geometric information it
encodes while insufficient, in general, to reconstruct the richer underlying
dynamics.

\paragraph{Our contribution.}

We first construct a canonical spectralization of finite differential data. A
star-compatible derivation $\delta$ into a real Hilbert $\A$-bimodule determines
a minimal one-form carrier and a real even Hodge--Dirac spectral triple with exact
order-zero and first-order relations. Two rigidity results identify what this
construction does not determine. The algebra and faithful trace do not fix an
absolute Dirac scale, and a fixed commutator derivation determines the full
compatible Dirac operator only modulo a finite-dimensional self-adjoint commutant.
The orthogonally centered representative is the unique
minimum-Hilbert--Schmidt-norm implementer, and the dimension of this commutant is
the sharp number of independent scalar measurements required to distinguish the
remaining Dirac ambiguity.

This distinction separates an \emph{austere Hodge construction} from an enlarged
\emph{calibrated marked reconstruction}. The bare Hodge operator always has the
unit zero-mode and therefore cannot realize arbitrary finite Dirac operators.
When the missing finite calibration data are retained, every finite spectral
packet is reconstructible and every separable compact-resolvent packet arises as
a strong cofinal limit of finite marked reconstructions. A stricter requirement,
in which multiplication must be transported asymptotically in norm, is genuinely
selective and is characterized exactly by spectral quasidiagonality.

We next test whether the construction has independent geometric content. For
stationary graph dynamics, the symmetric jump data determine a Hodge--Dirac
operator whose square is the symmetric Markov generator and whose commutator
seminorm is its local quadratic form. On periodic $A_3$ approximants, the
corresponding Connes distance converges uniformly to a Riemannian metric
reconstructed independently from the jump second moments. For closed Riemannian
spin manifolds, the local positive-stencil and covariant Wilson construction gives
the required chartwise estimates, while global norm-resolvent convergence follows
under the stated compatible stable-spin discretization hypotheses. The independent
scalar metric reconstruction is unconditional.

Finally, we characterize several mechanisms by which spectralization loses
information. Directed stationary current can vary at fixed metric spectral
geometry, hidden dynamical memory can vary at fixed static spectral data, and
locally identical spectral packets can differ by global marked monodromy. These
inverse fibres have intrinsic coordinates and canonical coarse-graining maps.
Thus spectralization is not merely noninjective: the information it forgets has
its own mathematical structure and transformation laws.

The resulting interpretation is that Connes-type spectral geometry can arise as
an effective geometric description of a richer predictive dynamics rather than
necessarily as its most primitive level. This conclusion is strengthened by the
companion spacetime construction: Lorentzian spacetime and spectral geometry can
both arise downstream of the same pregeometric predictive structure, without
either being required as the primitive entrance to the other. The two papers
therefore lead to the schematic picture
\begin{equation}
\boxed{
\begin{array}{c}
\text{operational histories}\\[1mm]
\downarrow\\[1mm]
\text{predictive Renewal Geometry}\\[1mm]
\downarrow\\[1mm]
\text{Renewal Grand Tensor}\\[1mm]
\swarrow\hspace{2.5cm}\searrow\\[-1mm]
\text{Lorentzian spacetime}
\hspace{1.4cm}
\text{Connes-type spectral geometry}.
\end{array}}
\label{eq:intro-renewal-geometric-branches}
\end{equation}

The left branch is developed in the companion Renewal Geometry
work~\cite{PelissierRenewalSpacetime2026}. The present paper develops the right
branch: its canonical finite construction, its category-dependent image, its
independent metric test, and its inverse fibres. The central question is therefore
not only why spectral geometry can arise, but what a richer underlying dynamics
retains that must be forgotten when it is compressed to a spectral description.

















% =============================================================================
\section{From predictive dynamics to spectral geometry}
\label{sec:finite-main}
% =============================================================================

The spectralization theorem is algebraically self-contained, but its motivating
entrance is more primitive than the datum $(\A,\tau,\delta)$.  We therefore begin
with the finite predictive construction of the Renewal Grand Tensor~\cite{PelissierRenewalSpacetime2026}, restated here in the finite form needed
for the present paper.  This makes explicit which structures are reconstructed
before spectral geometry is introduced and which first-order information enters
only at the spectralization stage.

\subsection{The Renewal Grand Tensor and the first-order differential interface}
\label{subsec:grand-tensor-main}

Fix a finite cutoff $X$.  The operational entrance consists of a finite typed
event grammar with finite reachable histories.  If a history $h$ ends at cut type
$x$, let $\mathsf F_{X,x}$ be the finite admissible future words beginning at $x$.
The conditional event law assigns normalized one-step probabilities and therefore
cylinder probabilities
\begin{equation}
 p_X(w\mid h)
 =\prod_{j=1}^{k}p_X(a_j\mid ha_1\cdots a_{j-1}),
 \qquad w=a_1\cdots a_k,
\label{eq:ncg-chain-law}
\end{equation}
with $p_X(\epsilon_x\mid h)=1$.  No persistent state space, algebra, Hilbert
space, metric, or Dirac operator is assumed at this stage.

Two histories of the same cut type are \emph{future equivalent} when no declared
finite continuation distinguishes them:
\begin{equation}
 h\sim_X h'
 \quad\Longleftrightarrow\quad
 p_X(w\mid h)=p_X(w\mid h')
 \quad\text{for every }w\in\mathsf F_{X,x}.
\label{eq:ncg-future-equivalence}
\end{equation}
Histories of different cut type are never identified.  This is a typed right
congruence.  Indeed, if $h\sim_Xh'$ and $a:x\to y$ is admissible, then the
one-letter future gives $p_X(a\mid h)=p_X(a\mid h')$; on a positive-probability
branch, for every continuation $v$,
\[
 p_X(v\mid ha)
 =\frac{p_X(av\mid h)}{p_X(a\mid h)}
 =\frac{p_X(av\mid h')}{p_X(a\mid h')}
 =p_X(v\mid h'a).
\]
Hence $ha\sim_Xh'a$.

Define the minimum persistent predictive carrier by
\begin{equation}
 \boxed{
 Z_{X,x}^{\min}:=\mathsf H_{X,x}/\!\sim_X,
 \qquad
 q_{X,x}:\mathsf H_{X,x}\twoheadrightarrow Z_{X,x}^{\min}.}
\label{eq:ncg-predictive-state}
\end{equation}
For an admissible event $a:x\to y$, right congruence gives canonical update and
branch maps
\begin{equation}
 \boxed{
 T_{X,a}([h])=[ha],
 \qquad
 \kappa_{X,a}([h])=p_X(a\mid h).}
\label{eq:ncg-predictive-update}
\end{equation}
Thus persistence is reconstructed from the conditional future law rather than
introduced as an independent coordinate.  Moreover, this quotient is universal:
if $S_X$ is any reachable deterministic predictive realization of the same
conditional event law, then equality of two realization states forces equality of
all their future cylinder probabilities.  Consequently there is a unique
surjection
\begin{equation}
 \pi_{X,x}:S_{X,x}\twoheadrightarrow Z_{X,x}^{\min}
\label{eq:ncg-predictive-universal}
\end{equation}
intertwining branch probabilities and updates.  Hence $Z_X^{\min}$ is the unique
coarsest reachable deterministic predictive realization.

To reconstruct the operator-algebraic history object, assume in addition that the
finite multitime intervention tables are specified on tomographically complete
finite operator frames, are affine under randomization and additive under coarse
graining, and obey the usual finite positivity, normalization, prefix
compatibility, and causal-discard conditions.  Finite tomography then determines
the corresponding positive causal process matrices.  Choosing their
support-minimal square-root realizations introduces no additional invariant:
minimal realizations of the same positive process data differ only by unitaries on
the support spaces.  The represented predictive-state projections, mixed forward
maps, and their Hilbert--Schmidt reverses therefore generate, up to represented
$*$-isomorphism, a finite history algebra
\begin{equation}
 \A_X^{\mathrm{hist}}\subseteq\mathcal B(\mathcal M_X).
\label{eq:ncg-history-algebra}
\end{equation}
Here $\mathcal M_X$ denotes any support-minimal sequential carrier; its particular
choice is immaterial up to the preceding unitary equivalence.

On every finite-dimensional $C^*$-algebra there is a canonical normalized regular
trace.  For
$\A_X^{\mathrm{hist}}\cong\bigoplus_\lambda M_{n_\lambda}(\C)$, set
\begin{equation}
 \widehat\tau_{X,\mathrm{reg}}(a)
 =\frac{\Tr_{\C}(L_a)}{\dim_{\C}\A_X^{\mathrm{hist}}}
 =\frac{\sum_\lambda n_\lambda\Tr(a_\lambda)}
        {\sum_\lambda n_\lambda^2}.
\label{eq:ncg-grand-regular-trace}
\end{equation}
If $[w]$ denotes the represented operator associated with a finite history word,
define the positive word-Gram kernels
\begin{equation}
 \mathbb K_{X,r}(v,w)
 =\widehat\tau_{X,\mathrm{reg}}\bigl([v]^*[w]\bigr),
 \qquad |v|,|w|\le r.
\label{eq:ncg-word-gram}
\end{equation}
Positivity is immediate from the trace inner product.  Since
$\A_X^{\mathrm{hist}}$ is finite dimensional, the increasing spans of represented
words stabilize after finite depth; once this happens, one additional word layer
determines multiplication and involution on the stabilized span.

\begin{definition}[Finite Renewal Grand Tensor]
\label{def:ncg-grand-tensor}
The finite Renewal Grand Tensor associated with the operational process $X$ is
\begin{equation}
 \boxed{
 \GTX
 =\left(
 \A_X^{\mathrm{hist}},
 \{\mathbb K_{X,r}\}_{r\ge0},
 Z_X^{\min}
 \right).}
\label{eq:ncg-grand-tensor}
\end{equation}
It is a tensor in the operational sense of a complete positive history-moment
object: $Z_X^{\min}$ records which past distinctions remain predictively effective,
$\A_X^{\mathrm{hist}}$ records how represented histories compose, and the Gram
hierarchy records their intrinsic positive overlaps.
\end{definition}

\begin{proposition}[Canonical finite Grand-Tensor reconstruction]
\label{prop:ncg-grand-tensor-reconstruction}
Under the finite assumptions above, the following data are fixed by the complete
operational statistics up to their natural isomorphisms:
\begin{enumerate}[label=\textup{(\roman*)},leftmargin=2.8em]
\item the future-minimal carrier $Z_X^{\min}$ and its update/branch maps;
\item the support-minimal represented history algebra
      $\A_X^{\mathrm{hist}}$, up to represented $*$-isomorphism;
\item the normalized regular trace and the positive word-Gram hierarchy;
\item consequently, the Grand Tensor $\GTX$ of
      \eqref{eq:ncg-grand-tensor}.
\end{enumerate}
No independent hidden-state presentation, persistent state space, metric, or
spectral triple is needed for this reconstruction.
\end{proposition}

\begin{proof}
Right congruence was proved above.  The universal factorization
\eqref{eq:ncg-predictive-universal} shows that $Z_X^{\min}$ is independent of a
chosen predictive presentation.  Tomographic completeness makes the finite
multitime positive process matrices unique; support-minimal Gram factorizations of
the same positive matrices are unitarily equivalent.  The represented operators
generated from such a realization therefore determine the same finite
$C^*$-algebra up to represented $*$-isomorphism.  The normalized regular trace is
intrinsic to that algebra, so the Gram kernels
\eqref{eq:ncg-word-gram} are invariant under the same equivalence.  Finite
dimensionality gives stabilization of the represented word spans, completing the
claim.
\end{proof}

This proposition is the self-contained predictive entrance needed in the present
paper.  The same construction is developed in the broader spacetime setting of the
Renewal Geometry framework~\cite{PelissierRenewalSpacetime2026}, including additional
same-history physical marks
needed by the Lorentzian branch; none of those spacetime-specific marks is assumed
for the spectralization below.

\paragraph{First-order differential interface.}

The Grand Tensor is logically upstream of spectral geometry, but it is important
not to identify algebraic reconstruction with metric differentiation.  From
\eqref{eq:ncg-grand-tensor} we canonically obtain the history algebra and its
positive finite-dimensional structure.  We take
\begin{equation}
 \A=\A_X^{\mathrm{hist}},
 \qquad
 \tau=\widehat\tau_{X,\mathrm{reg}},
\label{eq:ncg-grand-algebra-trace}
\end{equation}
or, more generally, any declared faithful tracial state on $\A$.  To spectralize
this algebra one additionally specifies a compatible first-order variation.
Abstractly, this is a Hilbert $\A$-bimodule $\mathcal K$, an antiunitary
involution $C_{\mathcal K}$ reversing the bimodule actions, and a complex-linear
map
\begin{equation}
 \delta:\A\longrightarrow\mathcal K
\label{eq:ncg-grand-delta}
\end{equation}
satisfying
\begin{equation}
 \delta(\one)=0,
 \qquad
 \delta(ab)=a\delta(b)+\delta(a)b,
 \qquad
 C_{\mathcal K}\delta(a)=\delta(a^*).
\label{eq:ncg-grand-delta-rules}
\end{equation}
When this datum comes from a Renewal process, $\delta$ represents the declared
first-order variation of the reconstructed predictive history algebra.  The
results below do not assume that $\GTX$ or $\A$ alone uniquely determines this
variation.  This distinction is essential: the no-scale result below shows that
rescaling $\delta$ changes the spectral metric while leaving the algebra and trace
fixed.

The logical interface is therefore
\begin{equation}
\boxed{
\text{operational histories}
\longrightarrow
\GTX
\longmapsto
(\A_X^{\mathrm{hist}},\tau_X)
\overset{\delta_X}{\longrightarrow}
(\A_X,\Hh_X,D_X,J_X,\gamma_X).}
\label{eq:ncg-grand-spectral-interface}
\end{equation}
The first arrows reconstruct predictive and algebraic structure; the final arrow
is the Hodge--Dirac spectralization studied in this paper.

The zero-form carrier is
\begin{equation}
\Hh_0=L^2(\A,\tau),
\label{eq:summary-H0}
\end{equation}
and the one-form carrier is the minimal quotient of the universal calculus
generated by $\delta$,
\begin{equation}
\Hh_1(\delta)=\Om_u^1(\A)/\Ker T_\delta,
\qquad
\Om_u^1(\A)=\Ker(m:\A\otimes\A\to\A).
\label{eq:summary-H1}
\end{equation}
On $\Hh=\Hh_0\oplus\Hh_1(\delta)$, the induced derivation $\partial$ gives
\begin{equation}
D=\begin{pmatrix}0&\partial^*\\ \partial&0\end{pmatrix},
\qquad
\gamma=\begin{pmatrix}I&0\\0&-I\end{pmatrix},
\label{eq:summary-D}
\end{equation}
while reversal of universal one-forms gives the real structure $J$.  Thus the
spectral triple is constructed from a first-order differential on an algebra that
has itself already been reconstructed from predictive histories.

\subsection{What the differential determines}

The finite Hodge construction is canonical once $\delta$ is fixed.  It produces a
real even spectral triple with exact order-zero and first-order relations.  The
spectral scale is not already hidden in $(\A,\tau)$: replacing $\delta$ by
$s\delta$ leaves the algebra and trace unchanged but sends
\begin{equation}
D\longmapsto sD,
\qquad
L_D\longmapsto sL_D,
\qquad
d_D\longmapsto s^{-1}d_D.
\label{eq:summary-scale}
\end{equation}
Thus an absolute metric scale enters through the differential data.

There is a second, finite ambiguity.  For a marked spectral packet
$\mathsf S=(\A,\Hh,\pi,D,J,\gamma)$ define
\begin{equation}
\mathfrak c_{\mathsf S}
=\{C\in\pi(\A)'_{\rm sa}: C\gamma+\gamma C=0,\ JC=\varepsilon'CJ\}.
\label{eq:summary-commutant}
\end{equation}
All compatible self-adjoint Dirac operators implementing the same commutator
derivation form exactly the affine space
\begin{equation}
D+\mathfrak c_{\mathsf S}.
\label{eq:summary-affine}
\end{equation}
The orthogonally centered representative
$D_{\rm der}=D-P_{\mathfrak c_{\mathsf S}}D$ is the unique
minimum-Hilbert--Schmidt-norm implementer, and
$q_{\mathsf S}=\dim_{\R}\mathfrak c_{\mathsf S}$ real linear scalar measurements are
necessary and sufficient to distinguish every point of the fibre.

The austere Hodge construction has an additional visible restriction:
$D(\one,0)=0$.  It therefore cannot be surjective onto arbitrary finite Dirac
operators.  This is why the inverse statement below uses an enlarged calibrated
marked reconstruction rather than the bare Hodge functor.

\subsection{Which spectral geometries can be reconstructed?}

At the finite marked level, one is allowed to retain finitely many calibration data,
including a target generator germ or the equivalent commutant coordinates.  At the
strong compact level, finite spectral compressions converge strongly on fixed
vectors, resolvents, commutators, and matrix coefficients.  At the strict
norm-monoidal level, multiplication must additionally be transported asymptotically
in norm.

\begin{theorem}[Spectralization, Dirac ambiguity, and calibrated image]
\label{thm:summary-finite}
The following statements hold.
\begin{enumerate}[label=\textup{(\roman*)},leftmargin=2.8em]
\item Every finite real differential datum $(\A,\tau,\delta)$ determines the real even
Hodge--Dirac spectral triple above, functorially under trace- and
 differential-preserving $*$-isomorphisms.
\item The algebra and trace do not determine the Dirac scale, and a fixed commutator
 derivation determines the full compatible Dirac operator exactly modulo
 $\mathfrak c_{\mathsf S}$.
\item For the enlarged calibrated marked reconstruction,
\begin{equation}
\begin{aligned}
\EssIm(\mathfrak S_{\rm fin}^{\rm mark})
  &=\mathbf{Spec}_{\rm fin},\\
\EssIm(\mathfrak S_{\rm cpt}^{\rm strong,mark})
  &=\mathbf{Spec}_{\rm cpt}^{\rm sep},\\
\EssIm(\mathfrak S_{\rm cpt}^{\rm strict})
  &=\mathbf{Spec}_{\rm cpt}^{\rm sqd}
  \subsetneq\mathbf{Spec}_{\rm cpt}^{\rm sep}.
\end{aligned}
\label{eq:summary-essential-image}
\end{equation}
Thus every finite spectral packet is reconstructible once the finite calibration
marks are retained; every separable compact-resolvent packet is a strong cofinal
limit of such finite marked reconstructions; and strict norm-monoidal realizability
is equivalent to spectral quasidiagonality.
\end{enumerate}
\end{theorem}

The distinction between the first and third clauses is essential.  Finite marked
surjectivity is a reconstruction benchmark for an enlarged data set; it is not a
claim that an arbitrary target $D$ has been derived from the algebra or Hodge
derivation alone.  The complete definitions and proofs are given in
\Cref{sec:finite} of the mathematical supplement.


% =============================================================================
\section{Independent metric reconstruction and compact spin geometry}
\label{sec:a3-main}
% =============================================================================

Producing an operator called $D$ is not by itself evidence that geometry has emerged.
A stronger test is available when the same finite dynamics determines a metric by a
route that does not use the Dirac operator.

\subsection{An independent metric test}

For a stationary finite graph, let $m_x>0$ be the stationary masses,
$q_{xy}=m_xk_{xy}$ the stationary fluxes, and
\begin{equation}
c_{xy}=\frac{q_{xy}+q_{yx}}2,
\qquad
j_{xy}=\frac{q_{xy}-q_{yx}}2.
\label{eq:summary-cj}
\end{equation}
The symmetric conductances determine the graph Hodge--Dirac operator $D_{m,c}$.
Exactly,
\begin{equation}
D_{m,c}^2|_{\ell^2(V,m)}=-2\mathcal L_s,
\qquad
\|[D_{m,c},f]\|^2
=\max_x\frac1{m_x}\sum_yc_{xy}|f(y)-f(x)|^2.
\label{eq:summary-graph-identities}
\end{equation}
Thus the Dirac square is the symmetric Markov generator and its commutator seminorm
is the local quadratic form of that generator.

On the periodic $A_3$ approximants, the normalization gives
\begin{equation}
L_h(f)^2
=\max_{x\in V_h}\frac1{8h^2}
\sum_{\alpha\in\Phi_{A_3}}|f(x+h\alpha)-f(x)|^2,
\label{eq:summary-A3-Lip}
\end{equation}
and the pure-state Connes distance satisfies
\begin{equation}
\sup_{x,y\in V_h}
|d_h^{\rm Connes}(x,y)-d_g(x,y)|\longrightarrow0,
\label{eq:summary-A3-distance}
\end{equation}
where $d_g$ is reconstructed independently from the jump second moments.  This is
the principal anti-circularity test of the paper.

\subsection{From local spin discretization to a global convergence criterion}

The spin problem is subtler.  Local positive stencils reconstruct the inverse metric,
and a density-symmetric covariant Wilson discretization removes lattice doublers while
remaining $O(h)$ on sampled smooth spinors.  These chartwise estimates do not by
themselves guarantee that independently discretized charts assemble into one stable
global operator.  The mathematical supplement therefore isolates the additional
global approximation properties in the notion of a \emph{compatible stable spin
discretization}.

\begin{theorem}[Independent metric reconstruction and stable compact-spin convergence]
\label{thm:summary-metric}
The following statements hold.
\begin{enumerate}[label=\textup{(\roman*)},leftmargin=2.8em]
\item The periodic $A_3$ Hodge--Dirac operator obeys
\eqref{eq:summary-graph-identities}, and its Connes distance converges uniformly as
in \eqref{eq:summary-A3-distance}.
\item The relation-adapted cellular completion removes the extensive one-form kernel
and has the expected compact Hodge--Dirac continuum limit.
\item The covariant Wilson construction has one low-energy species and satisfies the
local uniform graph estimates required for a spin continuum limit.
\item Let $(M^d,g,\mathfrak s)$ be closed.  If the finite covariant Wilson operators
form a compatible stable spin discretization in the precise sense of
Definition~\ref{def:stable-spin-atlas}, then for every $z\in\C\setminus\R$,
\begin{equation}
\left\|W_h(D_h-z)^{-1}W_h^*
-(\widehat D_{M,\mathfrak s}-z)^{-1}\right\|\longrightarrow0,
\qquad
\widehat D_{M,\mathfrak s}=D_{M,\mathfrak s}\otimes\sigma_1.
\label{eq:summary-spin-resolvent}
\end{equation}
\item Independently of the spin-assembly hypothesis, every closed Riemannian manifold
admits finite metric triples on refining meshes for which
\begin{equation}
\sup_{x,y\in V_h}|d_h^{\rm met}(x,y)-d_g(x,y)|\longrightarrow0.
\label{eq:summary-global-metric}
\end{equation}
\end{enumerate}
\end{theorem}

The theorem deliberately separates the verified local spin analysis from the global
atlas-stability requirement.  The scalar metric approximation is unconditional; the
spin norm-resolvent conclusion is a convergence criterion for compatible stable
assemblies.  The precise local estimates, atlas definition, and compactness argument
are given in \Cref{sec:metric,app:compact-spin}.


% =============================================================================
\section{Spectral geometry as a universality quotient}
\label{sec:universality-main}
% =============================================================================

The spectral map is broad but not injective.  Three independent mechanisms make the
loss of information explicit: directed current at fixed symmetric metric geometry,
hidden dynamical memory at fixed static spectral data, and global marked monodromy
at fixed local packets.

\subsection{Information lost under spectralization}

For a connected graph with fixed masses $m$ and conductances $c$, the complete
stationary current fibre is
\begin{equation}
\mathcal J(G,c)=\{j\in\Ker B:\ |j_e|<c_e\ \text{for every }e\},
\qquad
\dim\mathcal J(G,c)=b_1(G).
\label{eq:summary-current-fibre}
\end{equation}
Every point of this fibre has the same metric spectral triple, while path
probabilities, time-reversal asymmetry, and entropy production can vary.  The lost
directionality is recovered by the modular affinity
\begin{equation}
a_e=\operatorname{arctanh}\frac{j_e}{c_e}
=\frac12\log\frac{q_e}{q_{\bar e}},
\label{eq:summary-affinity}
\end{equation}
whose cohomology class parametrizes the current fibre.

Hidden positive memory gives a second fibre.  Eliminating auxiliary degrees of
freedom produces a Stieltjes self-energy
\begin{equation}
\Sigma_\mu(z)
=\int_{(0,\infty)}\frac{1}{\lambda-z}\,\dd\mu(\lambda)
=B_\mu(C_\mu-z)^{-1}B_\mu^*.
\label{eq:summary-Stieltjes}
\end{equation}
The full operator-valued measure $\mu$ classifies the minimal positive dilation,
whereas a static spectral head may see only selected boundary values or moments.
Finite rational Hermitian indefinite memory has the corresponding minimal
Pontryagin/descriptor classification.

The third fibre is global.  Locally marked-isomorphic spectral packets glued over a
connected incidence graph are classified by
\begin{equation}
[\rho]\in\Hom(\pi_1(\Gamma),G_{\mathsf S})/G_{\mathsf S}.
\label{eq:summary-monodromy}
\end{equation}
Trivial monodromy gives a global marked trivialization; nontrivial monodromy gives a
flat twisted spectral geometry.

\subsection{Universality fibres and coarse graining}

These fibres have canonical effective images.  Current is pushed through the graph
chain map; high-energy positive memory is replaced by an exact zero-energy contact
correction plus the retained low-energy measure; and monodromy descends precisely
when the collapsed cycles lie in its kernel.  The maps compose under successive
coarse grainings.

\begin{theorem}[Universality fibres and coarse graining]
\label{thm:summary-universality}
For the classes above:
\begin{enumerate}[label=\textup{(\roman*)},leftmargin=2.8em]
\item the stationary realization fibre of a metric graph triple is
$\mathcal J(G,c)$, and the modular class recovers the complete directed generator;
\item positive hidden memory is classified by its minimal operator-valued Stieltjes
measure, while finite rational Hermitian indefinite memory is classified by a
minimal Pontryagin/descriptor realization;
\item locally identical marked packets are globally classified by the monodromy class
\eqref{eq:summary-monodromy};
\item current, memory, and monodromy admit canonical composable coarse-graining maps,
with the corresponding entropy/rank/index monotonicity statements;
\item after auxiliary markings, calibration, and joint Gram data are retained, the
fully reduced finite inverse fibre is a stratified stabilizer quotient, and under
the explicit separation, tightness, domain, weight, and transport hypotheses of the
supplement the completed fibre is the inverse limit of those finite strata.
\end{enumerate}
\end{theorem}

The detailed current geometry, operator-measure realization, monodromy classification,
coarse-graining maps, rational indefinite memory, and completed inverse-fibre theorems
are given in \Cref{sec:inverse,app:rational-krein,app:completed-fibres}.




% =============================================================================
\section{Interpretation and discussion}
\label{sec:discussion}
% =============================================================================

The results separate predictive reconstruction from operator-theoretic
spectralization in a useful way.  Renewal Geometry first reconstructs the Grand
Tensor $\GTX$, including the future-minimal predictive carrier, represented
history algebra, and positive moment hierarchy.  Spectralization then equips that
reconstructed algebra with a compatible first-order differential datum
$\delta_X$.  Once $(\A,\tau,\delta)$ is given, the finite Hodge--Dirac theorem
is an ordinary statement about $C^*$-algebras, Hilbert bimodules, and derivations.
The predictive reconstruction is therefore logically prior to the spectral
triple without being an extra hypothesis of the abstract operator-theoretic
result.

\subsection{Spectral geometry as an effective description of predictive dynamics}

The first consequence is a change in the logical status of the spectral triple.
In Connes' formulation, the algebra, its Hilbert-space representation, and the
Dirac operator constitute the spectral data from which geometry is extracted.  In
the present construction these objects can instead occur downstream of the
Renewal Grand Tensor:
\begin{equation}
\text{predictive dynamics}
\longrightarrow
\GTX
\longmapsto
(\A_X^{\mathrm{hist}},\tau_X)
\overset{\delta_X}{\longrightarrow}
(\A_X,\Hh_X,D_X,J_X,\gamma_X).
\label{eq:discussion-layered}
\end{equation}
A spectral triple therefore need not be interpreted as the most microscopic
object in the description.  Within the framework considered here, it is the
spectral image of a specific algebraic and differential sector of a richer
predictive object.

This does not mean that an arbitrary Dirac operator is determined by the algebra.
The finite Hodge construction shows precisely where additional information enters.
The algebra and faithful trace do not fix an absolute metric scale, and the
commutator derivation determines the full compatible Dirac operator only modulo
the appropriate self-adjoint commutant.  Likewise, finite marked surjectivity is a
statement about an enlarged reconstruction category in which the missing
calibration coordinates are retained; it is not a claim that an arbitrary target
Dirac operator has been derived from the austere Hodge datum alone.  At the
stricter norm-monoidal level the reconstruction becomes genuinely selective, with
spectral quasidiagonality providing the exact criterion.

The independent metric construction is important for the same reason.  On the
periodic $A_3$ family, the Riemannian metric reconstructed from jump second moments
and the Connes distance reconstructed from the Hodge--Dirac operator are obtained
by different routes and converge to the same continuum geometry.  The agreement
therefore tests the spectralization rather than defining the target metric by
construction.  For general closed Riemannian manifolds the scalar metric
approximation remains unconditional, while the spin norm-resolvent statement
requires the compatible stable-discretization hypotheses stated above.

\paragraph{A common pregeometric origin of spacetime and spectral geometry.}

The interpretation is strengthened by the companion spacetime construction
\cite{PelissierRenewalSpacetime2026}.  There, the same Renewal Grand Tensor
$\GTX$ supports a reconstruction of Lorentzian spacetime without requiring a
spectral triple as an intermediate primitive.  Together, the two papers therefore
place the Grand Tensor as a common pregeometric object upstream of two distinct
downstream descriptions, as summarized in
\eqref{eq:intro-renewal-geometric-branches}.

This matters conceptually.  If spectral geometry were the only downstream output,
one could interpret the predictive construction merely as an alternative
procedure for generating spectral triples.  The existence of an independent
spacetime branch instead suggests that Lorentzian geometry and Connes-type
spectral geometry are different projections of a common pregeometric structure
rather than successive fundamental layers.

The claim is structural rather than ontological.  The results do not establish
that Renewal Geometry is the unique microscopic theory underlying all
noncommutative geometries or all physical spacetimes.  They show that, within the
framework studied here, neither spectral geometry nor spacetime geometry needs to
be placed at the primitive entrance.  Both may arise from a predictive structure
defined before either metric geometry or a spectral triple is introduced.

\paragraph{Information loss and nonfaithfulness.}

The inverse results sharpen this effective interpretation.  Spectralization is
generically many-to-one:
\begin{equation}
\mathfrak S(\mathcal R_1)=\mathfrak S(\mathcal R_2)
\qquad\text{need not imply}\qquad
\mathcal R_1\simeq\mathcal R_2,
\label{eq:discussion-nonfaithful}
\end{equation}
where $\mathcal R_i$ denote predictive or dynamical realizations and
$\mathfrak S$ their spectral image.  A spectral triple is therefore not, in
general, a complete invariant of the underlying dynamics.

The stationary-current fibre provides the simplest example.  Fixing the masses
and symmetric conductances fixes the metric Hodge--Dirac triple, while a
divergence-free stationary current remains undetermined.  Moving within this
fibre changes path probabilities, time-reversal asymmetry, and entropy production
without changing the spectral metric.  The spectral geometry retains the
symmetric metric sector while forgetting part of the directed dynamics.

The memory fibre exhibits a different form of information loss.  Eliminating
auxiliary degrees of freedom produces a frequency-dependent self-energy whose
minimal realization is encoded by an operator-valued measure.  A static spectral
description may retain only selected boundary values, moments, or an effective
operator, so distinct dynamical tails can have the same static spectral image.
What is lost here is temporal organization rather than merely a finite parameter.

Global marked monodromy provides a third mechanism.  Locally equivalent spectral
packets can differ in how their marked realizations are assembled globally.
Consequently spectralization can forget directional, dynamical, and global
information for mathematically different reasons.

These examples motivate the description of spectral geometry as an
\emph{information-reduced effective geometry}.  This should not be interpreted as
a claim that noncommutative geometry is mathematically incomplete.  A spectral
triple may completely encode the geometric observables for which it is designed.
The point is that it need not completely encode the predictive dynamics from which
it arose.  A single spectral geometry can therefore represent an entire family of
dynamically inequivalent microscopic realizations.

The analogy with effective descriptions in statistical physics is useful.  A set
of macroscopic variables may completely characterize an effective state while
failing to determine a unique microscopic realization.  Likewise, a spectral
triple may fully characterize an effective metric and differential sector while
identifying systems that remain distinct at the predictive level.

\paragraph{Why the quotient has a universality interpretation.}

The inverse fibre
\begin{equation}
\mathfrak S^{-1}(\mathsf S)
\label{eq:discussion-fibre}
\end{equation}
is therefore more than a technical obstruction to inversion.  It describes the
microscopic information quotiented out in the passage to a fixed spectral packet.

Moreover, the discarded variables are themselves organized under effective
reduction.  Directed currents are pushed through the induced graph chain map;
high-energy memory can be absorbed into a zero-energy contact correction while
the retained low-energy measure survives; and marked monodromy descends exactly
when the cycles removed by the quotient carry trivial holonomy.  These maps
compose under successive coarse grainings.

This is the sense in which the spectral packet acts as a universality-level
observable.  Different microscopic systems can share the same spectral geometry
because distinctions that remain meaningful at the predictive level no longer
affect the operator-geometric structure retained at that level of description.

The coarse-graining language should nevertheless be interpreted precisely.  The
results do not prove that spectralization is universally generated by a
renormalization-group flow, nor that arbitrary predictive dynamics converges
toward a spectral fixed point.  What is established is a controlled many-to-one
map, explicit inverse fibres, and composable effective maps on their principal
coordinates.  The universality interpretation is therefore structural rather than
a claim of a universal dynamical attractor.

\subsection{Inverse geometry, scope, and outlook}

Once the lost information has been identified, a natural question is whether it
should be retained when it is physically relevant.  The present examples suggest
several possible enrichments.  A metric spectral triple can be supplemented by a
modular or current class to retain irreversible directed information; an
operator-valued self-energy can retain frequency-dependent memory; and a marked
holonomy class can retain global assembly data.  Schematically,
\begin{equation}
(\A,\Hh,D,J,\gamma)
\quad\longrightarrow\quad
(\A,\Hh,D,J,\gamma;
\text{current},\text{memory},\text{monodromy},\ldots).
\label{eq:discussion-enrichment}
\end{equation}
The ordinary spectral packet is recovered by forgetting the additional
coordinates.

This perspective changes the natural inverse problem.  Rather than expecting one
microscopic dynamics to correspond to one spectral geometry, one should generally
expect a moduli space or stratified fibre of realizations.  The finite stabilizer
quotient and its completed inverse-limit version provide concrete examples of
such an inverse geometry.  The relevant microscopic question is therefore not
necessarily
\begin{quote}
``Which dynamics is represented by this spectral triple?''
\end{quote}
but rather
\begin{quote}
``Which class of predictive dynamics has this spectral triple as its common
effective geometry?''
\end{quote}

The companion spacetime construction makes this question still richer.  Two
predictive systems that have the same spectral image need not necessarily have
the same image under another downstream geometric reconstruction.  Conversely,
systems could agree on a spacetime sector while differing in spectral or
dynamical refinements.  The joint fibres of the spectral and Lorentzian
projections of Renewal Geometry therefore provide a natural object for future
study.

\paragraph{Mathematical scope.}

Several limitations of the present results should remain explicit.  The strict
spectral image would benefit from an intrinsic characterization independent of a
chosen spectral projection sequence.  For spin geometry, the local Wilson
construction provides the required analytic estimates, but an unconditional
global realization theorem would require verification of the compatibility and
stability hypotheses for every closed spin manifold.  Noncompact and genuinely
semifinite geometries require further control of domains, localization, and
weights.  The positive Stieltjes memory classification has a complete finite
rational indefinite analogue, but the unrestricted nonrational
infinite-dimensional indefinite problem is not treated here.

Nor does the finite compact-quantum-metric-space result by itself imply
quantum Gromov--Hausdorff, propinquity, or spectral-propinquity convergence in the
continuum limit \cite{Rieffel2004,Latremoliere2022,FarsiLatremolierePacker2024}.
Such statements require uniform control of the full state-space Lip-norms and, for
spectral propinquity, compatible control of the spectral data beyond the
pure-state distance convergence proved for the $A_3$ approximants.

\paragraph{Open questions and broader interpretation.}

The effective-geometric viewpoint raises several further questions.  Which parts
of the inverse fibre can be reconstructed from natural enrichments of a spectral
triple?  Which fibre coordinates remain stable under broader classes of
coarse graining?  Can additional dynamical principles select canonical
representatives inside a spectral fibre?  And how much of the many-to-one
structure persists in continuum classes substantially larger than the finite and
controlled compact-resolvent settings considered here?

The coexistence of the spectral and spacetime branches suggests a second family
of questions.  One may ask when the two downstream reconstructions determine one
another, when one retains information absent from the other, and which microscopic
coordinates are invisible to both.  More generally, one can ask whether other
geometric or physical structures arise as additional quotients of the same
predictive object.

Within the scope established in this paper, the hierarchy is nevertheless clear.
Predictive dynamics carries temporal, directional, and global information before
a metric or spectral triple is introduced.  Renewal Geometry organizes that
information in the pregeometric Grand Tensor $\GTX$.  Lorentzian spacetime can
emerge from one downstream sector, while on the spectral branch $\GTX$
canonically supplies the predictive history algebra and positive moment structure,
which become Connes-type spectral geometry after a compatible first-order
variation is specified.  The latter retains the metric and
differential structures relevant to its level while identifying classes of
underlying dynamics that remain distinct microscopically.

Thus, within this framework, Connes-type spectral geometry is naturally
interpreted as a stable effective geometry of a richer predictive dynamics:
powerful and potentially complete for its geometric sector, but not in general a
complete description of the dynamics from which that geometry emerges.






\begin{thebibliography}{99}
\bibitem{Connes1994}
A.~Connes,
\newblock \emph{Noncommutative Geometry},
\newblock Academic Press, San Diego, 1994.

\bibitem{RennieVarilly2008}
A.~Rennie and J.~C.~V\'arilly,
\newblock Reconstruction of manifolds in noncommutative geometry,
\newblock arXiv:math/0610418, revised 2008.

\bibitem{connesreconstruction}
A.~Connes,
\newblock On the spectral characterization of manifolds,
\newblock \emph{Journal of Noncommutative Geometry} \textbf{7} (2013), no.~1,
1--82.

\bibitem{Rieffel1999}
M.~A.~Rieffel,
\newblock Metrics on state spaces,
\newblock \emph{Documenta Mathematica} \textbf{4} (1999), 559--600.

\bibitem{Rieffel2004}
M.~A.~Rieffel,
\newblock Gromov--Hausdorff distance for quantum metric spaces,
\newblock \emph{Memoirs of the American Mathematical Society} \textbf{168}
(2004), no.~796, 1--65.
\newblock doi:10.1090/memo/0796.

\bibitem{Latremoliere2022}
F.~Latr\'emoli\`ere,
\newblock The Gromov--Hausdorff propinquity for metric spectral triples,
\newblock \emph{Advances in Mathematics} \textbf{404} (2022), 108393.
\newblock doi:10.1016/j.aim.2022.108393.

\bibitem{FarsiLatremolierePacker2024}
C.~Farsi, F.~Latr\'emoli\`ere, and J.~Packer,
\newblock Convergence of inductive sequences of spectral triples for the spectral
propinquity,
\newblock \emph{Advances in Mathematics} \textbf{437} (2024), 109442.
\newblock doi:10.1016/j.aim.2023.109442.

\bibitem{PelissierRenewalSpacetime2026}
A.~P\'elissier,
\newblock Renewal Geometry and the emergence of Lorentzian spacetime,
\newblock HAL preprint hal-05744772, 2026.
\newblock \url{https://hal.science/hal-05744772}

\bibitem{Krajewski1998}
T.~Krajewski,
\newblock Classification of finite spectral triples,
\newblock \emph{Journal of Geometry and Physics} \textbf{28} (1998), 1--30.
\newblock doi:10.1016/S0393-0440(97)00068-5.

\bibitem{ChristensenIvan2006}
E.~Christensen and C.~Ivan,
\newblock Spectral triples for AF $C^*$-algebras and metrics on the Cantor set,
\newblock \emph{Journal of Operator Theory} \textbf{56} (2006), no.~1, 17--46.

\bibitem{PearsonBellissard2009}
J.~Pearson and J.~Bellissard,
\newblock Noncommutative Riemannian geometry and diffusion on ultrametric Cantor sets,
\newblock \emph{Journal of Noncommutative Geometry} \textbf{3} (2009), no.~3,
447--480.

\bibitem{ChristensenIvanSchrohe2012}
E.~Christensen, C.~Ivan, and E.~Schrohe,
\newblock Spectral triples and the geometry of fractals,
\newblock \emph{Journal of Noncommutative Geometry} \textbf{6} (2012), no.~2,
249--274.
\newblock doi:10.4171/JNCG/91.

\bibitem{KellendonkSavinien2012}
J.~Kellendonk and J.~Savinien,
\newblock Spectral triples and characterization of aperiodic order,
\newblock \emph{Proceedings of the London Mathematical Society} \textbf{104} (2012),
no.~1, 123--157.

\bibitem{Paterson2014}
A.~L.~T.~Paterson,
\newblock Contractive spectral triples for crossed products,
\newblock \emph{Mathematica Scandinavica} \textbf{114} (2014), no.~2, 275--298.
\newblock doi:10.7146/math.scand.a-17112.

\bibitem{JulienPutnam2016}
A.~Julien and I.~F.~Putnam,
\newblock Spectral triples for subshifts,
\newblock \emph{Journal of Functional Analysis} \textbf{270} (2016), no.~3,
1031--1063.
\newblock doi:10.1016/j.jfa.2015.12.002.

\bibitem{AielloGuidoIsola2022}
V.~Aiello, D.~Guido, and T.~Isola,
\newblock Spectral triples on irreversible $C^*$-dynamical systems,
\newblock \emph{International Journal of Mathematics} \textbf{33} (2022), no.~1,
2250005.
\newblock doi:10.1142/S0129167X22500057.

\bibitem{Majid2025}
S.~Majid,
\newblock Dirac operator associated to a quantum metric,
\newblock \emph{Journal of Noncommutative Geometry} \textbf{19} (2025), no.~2,
573--600.

\bibitem{CiprianiSauvageot2003}
F.~Cipriani and J.-L.~Sauvageot,
\newblock Derivations as square roots of Dirichlet forms,
\newblock \emph{Journal of Functional Analysis} \textbf{201} (2003), no.~1,
78--120.

\bibitem{HinzTeplyaev2013}
M.~Hinz and A.~Teplyaev,
\newblock Dirac and magnetic Schr\"odinger operators on fractals,
\newblock \emph{Journal of Functional Analysis} \textbf{265} (2013), no.~11,
2830--2854.
\newblock doi:10.1016/j.jfa.2013.07.021.

\bibitem{HinzKelleherTeplyaev2015}
M.~Hinz, D.~J.~Kelleher, and A.~Teplyaev,
\newblock Metric and spectral triples for Dirichlet and resistance forms,
\newblock \emph{Journal of Noncommutative Geometry} \textbf{9} (2015), no.~2,
359--390.
\newblock doi:10.4171/JNCG/195.

\bibitem{Cipriani2016}
F.~Cipriani,
\newblock Noncommutative potential theory: A survey,
\newblock \emph{Journal of Geometry and Physics} \textbf{105} (2016), 25--59.

\bibitem{Arhancet2024}
C.~Arhancet,
\newblock Spectral triples, Coulhon--Varopoulos dimension and heat kernel estimates,
\newblock \emph{Advances in Mathematics} \textbf{451} (2024), 109794.

\bibitem{ChiribellaDArianoPerinotti2009}
G.~Chiribella, G.~M.~D'Ariano, and P.~Perinotti,
\newblock Theoretical framework for quantum networks,
\newblock \emph{Physical Review A} \textbf{80} (2009), 022339.

\bibitem{PollockProcessTensor2018}
F.~A.~Pollock, C.~Rodr\'iguez-Rosario, T.~Frauenheim, M.~Paternostro, and K.~Modi,
\newblock Non-Markovian quantum processes: Complete framework and efficient
characterization,
\newblock \emph{Physical Review A} \textbf{97} (2018), 012127.
\newblock doi:10.1103/PhysRevA.97.012127.

\bibitem{PollockEtAl2018}
F.~A.~Pollock, C.~Rodr\'iguez-Rosario, T.~Frauenheim, M.~Paternostro, and K.~Modi,
\newblock Operational Markov condition for quantum processes,
\newblock \emph{Physical Review Letters} \textbf{120} (2018), 040405.

\bibitem{Jaeger2000}
H.~Jaeger,
\newblock Observable operator models for discrete stochastic time series,
\newblock \emph{Neural Computation} \textbf{12} (2000), no.~6, 1371--1398.
\newblock doi:10.1162/089976600300015411.

\bibitem{LittmanSuttonSingh2001}
M.~L.~Littman, R.~S.~Sutton, and S.~Singh,
\newblock Predictive representations of state,
\newblock in \emph{Advances in Neural Information Processing Systems 14},
MIT Press, 2002, pp.~1555--1561.

\bibitem{SinghJamesRudary2004}
S.~Singh, M.~R.~James, and M.~R.~Rudary,
\newblock Predictive state representations: A new theory for modeling dynamical
systems,
\newblock in \emph{Proceedings of the 20th Conference on Uncertainty in Artificial
Intelligence}, 2004, pp.~512--519.

\bibitem{CrutchfieldYoung1989}
J.~P.~Crutchfield and K.~Young,
\newblock Inferring statistical complexity,
\newblock \emph{Physical Review Letters} \textbf{63} (1989), 105--108.
\newblock doi:10.1103/PhysRevLett.63.105.

\bibitem{shalizicrutchfield}
C.~R.~Shalizi and J.~P.~Crutchfield,
\newblock Computational mechanics: Pattern and prediction, structure and simplicity,
\newblock \emph{Journal of Statistical Physics} \textbf{104} (2001), 817--879.

\bibitem{TraversCrutchfield2025}
N.~F.~Travers and J.~P.~Crutchfield,
\newblock Equivalence of history and generator $\epsilon$-machines,
\newblock \emph{Symmetry} \textbf{17} (2025), no.~1, 78.
\newblock doi:10.3390/sym17010078.

\bibitem{MarzenCrutchfield2017}
S.~E.~Marzen and J.~P.~Crutchfield,
\newblock Structure and randomness of continuous-time, discrete-event processes,
\newblock \emph{Journal of Statistical Physics} \textbf{169} (2017), 303--315.

\bibitem{Rieffel1998}
M.~A.~Rieffel,
\newblock Metrics on states from actions of compact groups,
\newblock \emph{Documenta Mathematica} \textbf{3} (1998), 215--230.
\newblock doi:10.4171/DM/41.

\bibitem{lawson-michelsohn}
H.~B.~Lawson and M.-L.~Michelsohn,
\newblock \emph{Spin Geometry},
\newblock Princeton University Press, Princeton, 1989.

\bibitem{nielsenninomiya1981a}
H.~B.~Nielsen and M.~Ninomiya,
\newblock Absence of neutrinos on a lattice. I. Proof by homotopy theory,
\newblock \emph{Nuclear Physics B} \textbf{185} (1981), 20--40.

\bibitem{nielsenninomiya1981b}
H.~B.~Nielsen and M.~Ninomiya,
\newblock Absence of neutrinos on a lattice. II. Intuitive topological proof,
\newblock \emph{Nuclear Physics B} \textbf{193} (1981), 173--194.

\bibitem{Nakamura2024}
S.~Nakamura,
\newblock Remarks on discrete Dirac operators and their continuum limits,
\newblock \emph{Journal of Spectral Theory} \textbf{14} (2024), no.~1,
255--269.
\newblock doi:10.4171/JST/478.

\end{thebibliography}





\clearpage

% Manual link from the main TOC to the Mathematical Supplement
\phantomsection
\addcontentsline{toc}{section}{S.~Mathematical Supplement}

\mathsupplementtrue
\appendix
\renewcommand{\thesection}{S\arabic{section}}
\setcounter{section}{0}


\renewcommand{\thesection}{S\arabic{section}}
\setcounter{section}{0}

\begin{center}
{\LARGE\bfseries Mathematical Supplement}
\end{center}
\vspace{1em}

The main text reconstructs the finite Renewal Grand Tensor from operational
histories before introducing spectral data.  Its algebraic sector canonically
supplies the future-minimal predictive carrier, the represented history algebra,
and its positive moment structure, while a compatible first-order variation is
supplied separately.  The mathematical supplement begins at this algebraic
interface, since the spectralization theorems depend only on the resulting
finite-dimensional unital $C^*$-algebra, faithful tracial state, and compatible
derivation into a Hilbert bimodule.

The associated Hodge--Dirac construction yields a real even spectral triple,
whose metric scale is not fixed by the algebra and trace alone, and whose Dirac
operator is determined by its commutator derivation only modulo a compatible
self-adjoint commutant.  After adjoining finite calibrated markings, every finite
spectral packet is reconstructible; every separable compact-resolvent packet arises
as a strong cofinal limit of such finite reconstructions, while norm-monoidal
realizability is equivalent to spectral quasidiagonality.  For stationary Markov
graphs, the Hodge--Dirac square is the symmetric generator and its commutator
seminorm is the corresponding local quadratic form; on periodic $A_3$
approximants, the resulting Connes distance converges uniformly to the Riemannian
metric reconstructed independently from the jump second moments.  The same
framework also makes the nonfaithfulness of spectralization explicit through
stationary currents, dynamical memory, and global marked monodromy, together with
their natural coarse-graining maps.  The sections below give the precise
definitions, hypotheses, constructions, and proofs underlying these results.

\medskip
\supplementtableofcontents

% =============================================================================
\section{Finite spectralization and essential image}
\label{sec:finite}
% =============================================================================

\subsection{Algebraic input and operational motivation}
\label{subsec:algebraic-input}

The finite construction requires only a finite-dimensional unital $C^*$-algebra
$\A$, a faithful tracial state $\tau$, and the differential datum specified
below.  No operational reconstruction theorem is assumed.  The Hilbert space
of zero-forms is the vector space $\A$ with inner product
\begin{equation}
 \Hh_0=L^2(\A,\tau),\qquad
 \langle a,b\rangle_0=\tau(a^*b).
 \label{eq:history-H0}
\end{equation}
Faithfulness makes this inner product nondegenerate.  Left and right
multiplication are bounded commuting actions, with adjoints $L_a^*=L_{a^*}$
and $R_a^*=R_{a^*}$.  The right action is a representation of the opposite
algebra: $R_aR_b=R_{ba}$.

There is a distinguished possible choice of trace, but the results do not
require it.  If $\A\cong\bigoplus_\lambda M_{n_\lambda}(\C)$, its normalized
regular trace is
\begin{equation}
 \widehat\tau_{\rm reg}(a)
 =\frac{\Tr_{\C}(L_a)}{\dim_{\C}\A}
 =\frac{\sum_\lambda n_\lambda\Tr(a_\lambda)}
        {\sum_\lambda n_\lambda^2}.
 \label{eq:regular-trace-selfcontained}
\end{equation}
This is a faithful tracial state invariant under $*$-isomorphisms.  Thus both
$\A$ and $\widehat\tau_{\rm reg}$ have ordinary finite-dimensional algebraic
meanings independent of the motivating history construction.

A Hilbert $\A$-bimodule means a Hilbert space $\mathcal K$ with commuting
unital $*$-representations of $\A$ and $\A^{\rm op}$.  We denote these by
$\lambda_{\mathcal K}(a)$ and $\rho_{\mathcal K}(b)$, so that
$a\xi b=\lambda_{\mathcal K}(a)\rho_{\mathcal K}(b)\xi$ and
$\rho_{\mathcal K}(a)\rho_{\mathcal K}(b)=\rho_{\mathcal K}(ba)$.

\begin{definition}[Finite real differential datum]
\label{def:finite-real-differential}
A finite real differential datum on $(\A,\tau)$ consists of a Hilbert
$\A$-bimodule $\mathcal K$, an antiunitary involution $C_{\mathcal K}$ with
\[
 C_{\mathcal K}(a\xi b)=b^*C_{\mathcal K}(\xi)a^*,
\]
and a complex-linear map $\delta:\A\to\mathcal K$ satisfying
\begin{equation}
 \delta(\one)=0,\qquad
 \delta(ab)=a\delta(b)+\delta(a)b,\qquad
 C_{\mathcal K}\delta(a)=\delta(a^*).
 \label{eq:predictive-derivation}
\end{equation}
Only the bimodule generated by $\delta(\A)$ is retained.  The term
\emph{predictive differential datum} will occasionally be used when this same
algebraic datum is obtained from the operational setting.  No result in this
section depends on how the datum was produced.
\end{definition}

For completeness, the one-form construction can be performed entirely within
this datum.  The universal first-order calculus is
\begin{equation}
 \Om_u^1(\A)=\Ker(m:\A\otimes\A\to\A),\qquad
 \dd_u a=\one\otimes a-a\otimes\one,
 \label{eq:universal-oneforms}
\end{equation}
with bimodule action $c(a\otimes b)d=ca\otimes bd$ and conjugate-linear reversal
\begin{equation}
 \left(\sum_i a_i\otimes b_i\right)^\natural
 =-\sum_i b_i^*\otimes a_i^*.
 \label{eq:oneform-reversal}
\end{equation}
Direct expansion gives $(\dd_u a)^\natural=\dd_u(a^*)$ and
$\dd_u(ab)=a\dd_u b+(\dd_u a)b$.
The map
\[
 T_\delta\left(\sum_i a_i\otimes b_i\right)
 =\sum_i a_i\delta(b_i),\qquad \sum_i a_ib_i=0,
\]
is the unique bimodule map with $T_\delta\dd_u=\delta$.  Indeed, the
Leibniz rule makes it right linear because the additional term is
$(\sum_i a_ib_i)\delta(c)=0$; left linearity is immediate.
Star compatibility gives
$T_\delta(\omega^\natural)=C_{\mathcal K}T_\delta(\omega)$.
Consequently
\begin{equation}
 Q_\delta(\omega,\eta)
 =\langle T_\delta\omega,T_\delta\eta\rangle_{\mathcal K}
 \label{eq:differential-Gram}
\end{equation}
is positive, and its null space is exactly $\Ker T_\delta$.  The resulting
one-form Hilbert space is
\begin{equation}
 \Hh_1(\delta)=\Om_u^1(\A)/\Ker T_\delta
 \cong\Ran T_\delta
 =\Span\{a\delta(b)c:a,b,c\in\A\}.
 \label{eq:H1-delta}
\end{equation}
It is finite dimensional even if the initially supplied $\mathcal K$ is not.
The displayed identification is isometric, so this quotient removes exactly
the unused bimodule directions.  This is the meaning of \emph{source minimality}
here.  Write $\partial a=[\dd_u a]$ for the induced derivation.

Predictive histories provide a motivation and a class of examples, rather than
extra axioms for this construction.  Section~\ref{subsec:grand-tensor-main}
reconstructs self-containedly the finite Renewal Grand Tensor $\GTX$ from
operational histories.  In particular, $\GTX$ canonically contains the separated
history algebra $\A_X^{\mathrm{hist}}$ and its positive word-moment hierarchy;
the normalized regular trace is an intrinsic faithful tracial state on that
finite algebra.  A compatible first-order variation $\delta_X$ is then supplied
as the additional differential datum to be spectralized, as described in
Section~\ref{subsec:grand-tensor-main}.  Once $(\A,\tau,\delta)$ is given,
the finite results below use only the algebraic data of
Definition~\ref{def:finite-real-differential}.

\subsection{The associated spectral triple}

Set $\Hh_0=L^2(\A,\tau)$ and let $L_a$ and $R_a$ denote left and right
multiplication on $\Hh_0$.  On
\begin{equation}
 \Hh=\Hh_0\oplus\Hh_1(\delta)
 \label{eq:H-total}
\end{equation}
define
\begin{equation}
 \pi(a)=L_a\oplus\lambda(a),
 \qquad
 D=\begin{pmatrix}0&\partial^*\\ \partial&0\end{pmatrix},
 \qquad
 \gamma=\begin{pmatrix}I&0\\0&-I\end{pmatrix}.
 \label{eq:finite-D}
\end{equation}
On $\Hh_0$ put $J_0a=a^*$, and on $\Hh_1(\delta)$ let
$J_1[\omega]=[\omega^\natural]$; set $J=J_0\oplus J_1$.

\begin{theorem}[Finite spectralization]
\label{thm:finite-spectralization}
The quintuple $(\A,\Hh,D,J,\gamma)$ is a real even finite spectral triple.
More precisely, $\pi$ is faithful, $D$ is self-adjoint, $J$ is antiunitary,
\[
 J^2=I,\qquad JD=DJ,\qquad J\gamma=\gamma J,
\]
and the opposite representation is
\[
 \pi^o(b)=J\pi(b^*)J^{-1}=R_b\oplus\rho(b).
\]
For all $a,b\in\A$,
\begin{equation}
 [\pi(a),\pi^o(b)]=0,
 \qquad
 [[D,\pi(a)],\pi^o(b)]=0.
 \label{eq:order-zero-first}
\end{equation}
If $B_a:\Hh_0\to\Hh_1(\delta)$ is defined by
$B_ab=(\partial a)b$, then
\begin{equation}
 [D,\pi(a)]
 =\begin{pmatrix}0&-B_{a^*}^*\\B_a&0\end{pmatrix}.
 \label{eq:finite-commutator}
\end{equation}
In particular, for $a=a^*$,
\begin{equation}
 \|[D,\pi(a)]\|=\|B_a\|,
 \qquad
 \Ker\bigl(a\mapsto\|[D,\pi(a)]\|\bigr)=\Ker\partial\cap\A_{\rm sa}.
 \label{eq:finite-Lip}
\end{equation}
\end{theorem}

\begin{proof}
Faithfulness follows from the left regular summand.  Since all spaces are finite
dimensional, $D$ is self-adjoint with compact resolvent.  The Leibniz rule gives
\[
 (\partial L_a-\lambda(a)\partial)b
 =\partial(ab)-a\partial b=(\partial a)b=B_ab,
\]
and taking adjoints yields \eqref{eq:finite-commutator}.  The reversal identities
make $J_1$ antiunitary and imply $J_1\partial=\partial J_0$, hence $JD=DJ$.
Moreover $J_0L_{b^*}J_0=R_b$ and
$J_1\lambda(b^*)J_1=\rho(b)$, so the opposite representation has the asserted
form.  Left and right actions commute.  Finally, the lower-left block of the
double commutator is $B_aR_b-\rho(b)B_a$, which vanishes because $B_a$ is right
$\A$-linear.  The upper-right block is its adjoint.  For self-adjoint $a$, the two
off-diagonal blocks are adjoints, proving \eqref{eq:finite-Lip}.
\end{proof}

\begin{remark}
The real structure in Theorem~\ref{thm:finite-spectralization} is the algebraic
reversal associated with the history bimodule.  Additional physical interpretations
of $J$ require extra markings and are not used in the mathematical results below.
\end{remark}

\begin{proposition}[No algebra-and-trace determination of the Dirac scale]
\label{prop:no-scale}
Let $\delta\neq0$ be a finite real differential datum and let $s>0$.  Replacing
$\delta$ by $s\delta$ leaves $(\A,\tau)$ unchanged and, under the canonical
identification of the one-form carriers, sends
\begin{equation}
 D\longmapsto sD,
 \qquad
 L_D(a):=\|[D,\pi(a)]\|\longmapsto sL_D(a).
 \label{eq:scale-D}
\end{equation}
Whenever the Connes distance is finite, it scales as
$d_D\mapsto s^{-1}d_D$.  Hence the algebra and canonical trace do not determine an
absolute spectral metric scale.
\end{proposition}

\begin{proof}
The inner product induced by $s\delta$ is $s^2Q_\delta$, so the canonical unitary
from the scaled one-form carrier to the original one sends the derivation vector
$[\dd_u a]$ to $s[\dd_u a]$.  Both off-diagonal blocks of $D$ therefore acquire the
factor $s$.  The commutator and dual-distance scalings follow immediately.
\end{proof}

The spectral seminorm also gives a compact quantum metric space, provided that
zero differential means constant observable.  This is the finite-dimensional
form of the state-space criterion underlying Rieffel's theory
\cite{Rieffel1998,Rieffel1999}.

\begin{proposition}[Compact quantum metric-space structure]
\label{prop:finite-CQMS}
Suppose $\Ker\partial=\C\one$ and set
$L_D(a)=\|[D,\pi(a)]\|$.  On the real order-unit space $\A_{\rm sa}$,
$L_D$ is a continuous Lip-norm.  Explicitly, the distance
\begin{equation}
 \operatorname{mk}_{L_D}(\varphi,\psi)
 =\sup\{|\varphi(a)-\psi(a)|:
          a=a^*,\ L_D(a)\le1\}
 \label{eq:finite-state-metric}
\end{equation}
is finite on the state space $S(\A)$ and induces its weak-$*$ topology.
Moreover, the complex extension satisfies
\begin{equation}
 L_D(ab)\le\|a\|L_D(b)+\|b\|L_D(a),\qquad
 L_D(a^*)=L_D(a).
 \label{eq:Leibniz-Lip}
\end{equation}
Conversely, if $\Ker\partial\ne\C\one$, the distance in
\eqref{eq:finite-state-metric} is infinite for some pair of states, so it is
not a compact quantum metric on the whole state space.
\end{proposition}

\begin{proof}
The commutator is a linear map between finite-dimensional normed spaces;
hence $L_D$ is continuous.  The commutator product rule and the $*$-identity
give \eqref{eq:Leibniz-Lip}.  By Theorem~\ref{thm:finite-spectralization},
the kernel on $\A_{\rm sa}$ is $\R\one$.
On $E_\tau=\{a=a^*:\tau(a)=0\}$ it follows that $L_D$ is a norm, and
finite-dimensional norm equivalence gives a constant $C_\delta$ such that
\begin{equation}
 \|a-\tau(a)\one\|\le C_\delta L_D(a),\qquad a=a^*.
 \label{eq:finite-Poincare-Lip}
\end{equation}
Thus $\{a\in E_\tau:L_D(a)\le1\}$ is norm compact.  Scalar multiples of
$\one$ do not affect differences of states, so the supremum in
\eqref{eq:finite-state-metric} can be taken over this compact set and is at
most $2C_\delta$.  Weak-$*$ convergence of states is uniform on a norm-compact
set, since states have norm one; it therefore implies convergence in
$\operatorname{mk}_{L_D}$.  Conversely, for each self-adjoint $a$,
\[
 |\varphi(a)-\psi(a)|\le L_D(a)
       \operatorname{mk}_{L_D}(\varphi,\psi),
\]
with the scalar case immediate.  Hence metric convergence implies weak-$*$
convergence and the metric separates states.

If the kernel is larger, star compatibility gives a non-scalar self-adjoint
$b\in\Ker\partial$.  States separate self-adjoint elements modulo scalars,
so some $\varphi,\psi$ have $\varphi(b)\ne\psi(b)$.  Every multiple $tb$
has zero seminorm, making their distance infinite.
\end{proof}

The quotient used here is $\A_{\rm sa}/\R\one$; the metric itself is on
$S(\A)$, not on a quotient of states.  Finite dimensionality is essential to
the automatic compactness argument: for an infinite-dimensional limit, a scalar
kernel alone does not replace total boundedness of the normalized Lipschitz ball.


\begin{proposition}[Unit zero-mode of the austere Hodge construction]
\label{prop:austere-zero-mode}
For every differential datum of Definition~\ref{def:finite-real-differential}, the
Hodge operator of Theorem~\ref{thm:finite-spectralization} satisfies
\begin{equation}
 D(\one,0)=0.
 \label{eq:austere-zero-mode}
\end{equation}
Consequently the bare Hodge functor $\mathfrak H_{\rm fin}$ never produces an
invertible finite Dirac operator.  More generally, knowledge of the commutator
derivation does not determine the component of $D$ in the compatible commutant;
that residual ambiguity is characterized in Theorem~\ref{thm:commutant-fibre}.
\end{proposition}

\begin{proof}
The universal derivation satisfies $\partial\one=0$, hence the lower-left block of
$D(\one,0)$ vanishes; the upper-right block acts on the zero one-form component and
vanishes as well.  The second assertion is immediate from
Theorem~\ref{thm:commutant-fibre}.
\end{proof}

\subsection{The commutant fibre of a fixed derivation}

Let $\mathsf S=(\A,\Hh,\pi,D,J,\gamma)$ be a finite marked spectral packet.  Define
\begin{equation}
 \mathfrak c_{\mathsf S}
 =\{C\in\pi(\A)'_{\rm sa}:
 C\gamma+\gamma C=0,\quad JC=\varepsilon'CJ\},
 \label{eq:commutant-space}
\end{equation}
with the evident omission of a condition when the corresponding mark is absent.

\begin{theorem}[Exact Dirac commutant fibre]
\label{thm:commutant-fibre}
Among self-adjoint operators with the declared grading and reality signs that
implement the same commutator derivation $[D,\pi(\cdot)]$, the complete solution
set is the affine space
\begin{equation}
 D+\mathfrak c_{\mathsf S}.
 \label{eq:Dirac-affine}
\end{equation}
Let $P_{\mathfrak c_{\mathsf S}}$ denote Hilbert--Schmidt projection onto
$\mathfrak c_{\mathsf S}$ and put
\begin{equation}
 D_{\rm der}=D-P_{\mathfrak c_{\mathsf S}}D.
 \label{eq:centered-D}
\end{equation}
Then $D_{\rm der}$ is the unique minimum-Hilbert--Schmidt-norm representative of
the affine fibre.  Moreover
$q_{\mathsf S}=\dim_{\R}\mathfrak c_{\mathsf S}$ independent real linear scalar measurements
are necessary and sufficient to distinguish every point of the fibre.
\end{theorem}

\begin{proof}
If $D'$ has the same commutator derivation as $D$, then
$C=D'-D$ commutes with $\pi(\A)$.  Subtracting the grading and reality equations
for $D'$ and $D$ shows that $C\in\mathfrak c_{\mathsf S}$.  Conversely every
$C\in\mathfrak c_{\mathsf S}$ preserves the commutators and the declared signs.
Thus the solution set is \eqref{eq:Dirac-affine}.  Orthogonal projection gives
\[
 \|D+C\|_{\HS}^2
 =\|D_{\rm der}\|_{\HS}^2
  +\|P_{\mathfrak c_{\mathsf S}}D+C\|_{\HS}^2,
\]
which proves the minimum statement.  Finally, a bank of scalar anchors restricts
to a real linear map on $\mathfrak c_{\mathsf S}$; injectivity requires at least
$q_{\mathsf S}$ coordinates and is achieved by any basis of the dual space.
\end{proof}

\subsection{Functoriality and essential image}

Let $\mathbf{RenDiff}_{\rm fin}$ have objects $(\A,\tau,\delta)$ as above and
morphisms the trace-preserving $*$-isomorphisms intertwining the differential
bimodule data.  The construction of Theorem~\ref{thm:finite-spectralization} is
functorial: the algebra isomorphism induces unitaries on $L^2(\A,\tau)$ and on the
minimal one-form quotient, and these unitaries intertwine $\pi,D,J,\gamma$.
We denote this austere Hodge functor by $\mathfrak H_{\rm fin}$.

The marked reconstruction considered next is a different, deliberately enlarged
map.  It retains a represented carrier together with finitely many calibrated
operator marks in addition to the differential datum; we denote it by
$\mathfrak S_{\rm fin}^{\rm mark}$.  The distinction is essential.  The Hodge
operator of Theorem~\ref{thm:finite-spectralization} annihilates $(\one,0)$, so the
bare functor $\mathfrak H_{\rm fin}$ cannot be unitarily equivalent to an arbitrary
invertible finite Dirac operator.  The marked image statement below concerns the
enlarged calibrated reconstruction, not surjectivity of the Hodge functor.

For the inverse problem we distinguish three categories.  In the \emph{finite
marked category}, a finite target may retain finitely many calibration and marking
data.  In the \emph{strong compact category}, finite-rank spectral compressions converge
strongly on fixed vectors, resolvents, commutators, and matrix coefficients.  In
the \emph{strict norm-monoidal category}, the finite compressions must additionally
transport multiplication asymptotically in norm.  A compact-resolvent packet is
called \emph{spectrally quasidiagonal} if there exist finite-rank projections
$P_n\uparrow I$ with
\begin{equation}
 P_nD=DP_n,
 \qquad
 \|[P_n,\pi(a)]\|\longrightarrow0
 \label{eq:sqd}
\end{equation}
for every $a$ in a dense spectral algebra, equivariantly for $J$ and $\gamma$ when
these are marked.

\begin{theorem}[Calibrated marked image and norm-compatible essential image]
\label{thm:essential-image}
The following statements concern the enlarged calibrated reconstruction in the
first two lines and the norm-compatible approximation problem in the third.  They
do not assert that the austere Hodge functor $\mathfrak H_{\rm fin}$ is surjective:
\begin{equation}
 \begin{aligned}
 \EssIm(\mathfrak S_{\rm fin}^{\rm mark})
   &=\mathbf{Spec}_{\rm fin},\\
 \EssIm(\mathfrak S_{\rm cpt}^{\rm strong,mark})
   &=\mathbf{Spec}_{\rm cpt}^{\rm sep},\\
 \EssIm(\mathfrak S_{\rm cpt}^{\rm strict})
   &=\mathbf{Spec}_{\rm cpt}^{\rm sqd}
   \subsetneq\mathbf{Spec}_{\rm cpt}^{\rm sep}.
 \end{aligned}
 \label{eq:essential-image}
\end{equation}
Thus every finite spectral packet is reconstructible once the finite calibration
marks defining the marked category are retained; every separable compact-resolvent
packet is a strong cofinal limit of such finite marked reconstructions; and strict
norm-monoidal realizability is equivalent to spectral quasidiagonality.
\end{theorem}

\begin{proof}
For a finite target $(\mathcal B,K,\pi_0,D_0,J_0,\gamma_0)$, choose finitely
many self-adjoint generators of $\pi_0(\mathcal B)$.  Their exponentials, with
angles chosen so that functional calculus recovers the generators, give a finite
family of generating unitary matrices.  Retain their full matrix-word
coefficients on $K$, the declared $J_0,\gamma_0$ marks, and the calibrated
unitary germ $U(t)=e^{-itD_0}$.  At this point $D_0$ is deliberately part of the
calibration data: the identity $i\dot U(0)=D_0$ is a reconstruction of the marked
generator, not a derivation of $D_0$ from $\mathcal B$ or from a Hodge differential.
If only the conjugation germ is retained, its derivative determines the commutator
action instead, and the missing compatible commutant coordinates are exactly those
of Theorem~\ref{thm:commutant-fibre}.  Evaluation of these finite-dimensional
marked data recovers the target up to unitary isomorphism.  The first equality in
\eqref{eq:essential-image} should therefore be read as a calibrated reconstruction
benchmark: it identifies which finite information closes the gap from the austere
derivation image to an arbitrary marked target.  The predictive entrance and its Grand-Tensor interface are reconstructed
self-containedly in Section~\ref{subsec:grand-tensor-main}; the companion Renewal
Geometry work~\cite{PelissierRenewalSpacetime2026} develops the same predictive
construction in the broader Lorentzian-spacetime setting.

Let now $(\A,\Hh,\pi,D)$ be separable with compact resolvent.  Let $P_n$ be the
increasing finite spectral projections of $D$.  Then $P_nD=DP_n$ and $P_n\to I$
strongly.  For every fixed $a,b$,
\begin{equation}
 P_n\pi(ab)P_n-P_n\pi(a)P_n\pi(b)P_n
 =P_n\pi(a)(I-P_n)\pi(b)P_n
 \label{eq:compression-defect}
\end{equation}
tends strongly to zero.  Here the finite-stage observables are the compressed matrices
$P_n\pi(a)P_n$; the compression map from the original algebra is not asserted
to be a $*$-homomorphism at finite $n$.  This is precisely why its product defect
is retained.  Applying the finite marked construction to these matrix data
and $P_nDP_n$ yields strong convergence of fixed resolvents, commutators, and
matrix coefficients, proving the second equality.

If \eqref{eq:sqd} holds, \eqref{eq:compression-defect} tends to zero in norm, so
the spectral packet lies in the strict image.  Conversely, norm multiplicativity
applied to $(a^*,a)$ gives
\[
 \|P_n\pi(a)^*(I-P_n)\pi(a)P_n\|\to0,
\]
which is exactly
$\|(I-P_n)\pi(a)P_n\|^2\to0$.  Applying the same argument to $a^*$ controls the
opposite off-diagonal block, hence $\|[P_n,\pi(a)]\|\to0$.  This proves equality
with the spectrally quasidiagonal class.

Properness is seen from the Toeplitz representation on $\ell^2(\N_0)$ with number
operator $Ne_k=ke_k$ and unilateral shift $Se_k=e_{k+1}$.  Any finite-rank
projection commuting with $N$ is diagonal in the number basis, and every nontrivial
finite diagonal projection has $\|[P,S]\|=1$.  Hence this compact-resolvent spectral
packet lies in the strong image but not in the strict one.
\end{proof}

\begin{remark}
The first equality in \eqref{eq:essential-image} is intentionally not an
``emergence from the algebra'' statement.  The full generator germ is allowed as a
calibration mark, so finite surjectivity by itself is formal.  The nontrivial
rigidity content is instead the comparison with the austere Hodge image:
Proposition~\ref{prop:austere-zero-mode} gives its unavoidable unit zero-mode, and
Theorem~\ref{thm:commutant-fibre} identifies exactly what a fixed commutator
derivation fails to determine.  The strict norm-monoidal line adds a separate,
genuine approximation constraint.
\end{remark}


% =============================================================================
\section{Metric realization and compact spin geometry}
\label{sec:metric}
% =============================================================================

The essential-image theorem is a structural existence result.  To test whether the
spectralization also has geometric content, one needs a metric reconstructed by a
route that does not use the Dirac operator.  Stationary graph dynamics provides a
particularly transparent setting because the symmetric jump covariance determines
such an independent metric quantity.

\subsection{Stationary graphs and the Hodge--Dirac operator}

\begin{definition}[Stationary finite renewal graph]
\label{def:supp-renewal-graph}
Let $V$ be finite, let $m_x>0$ with $\sum_xm_x=1$, and let $k_{xy}\ge0$ for
$x\ne y$.  Put
\begin{equation}
 (\mathcal Lf)(x)=\sum_{y\ne x}k_{xy}(f(y)-f(x)),
 \qquad q_{xy}=m_xk_{xy},
 \label{eq:supp-renewal-generator}
\end{equation}
and assume stationarity
\begin{equation}
 \sum_y(q_{xy}-q_{yx})=0
 \qquad(x\in V).
 \label{eq:supp-stationarity}
\end{equation}
Define
\begin{equation}
 c_{xy}=\frac{q_{xy}+q_{yx}}2=c_{yx},
 \qquad
 j_{xy}=\frac{q_{xy}-q_{yx}}2=-j_{yx}.
 \label{eq:supp-conductance-current}
\end{equation}
Assume the undirected graph on the pairs with $c_{xy}>0$ is connected.
\end{definition}

Set
\begin{equation}
 \Hh_0=\ell^2(V,m),
 \qquad
 \Hh_1=\ell^2(E^{\rm or},c),
 \qquad
 (\partial f)(x,y)=f(y)-f(x),
 \label{eq:supp-graph-carriers}
\end{equation}
where both orientations are retained and
$\norm{\omega}_1^2=\sum_{x,y}c_{xy}|\omega(x,y)|^2$.  Let
\begin{equation}
 (\lambda(a)\omega)(x,y)=a(y)\omega(x,y),
 \qquad
 (\rho(a)\omega)(x,y)=a(x)\omega(x,y).
 \label{eq:supp-graph-actions}
\end{equation}
Then $\partial(fg)=\lambda(f)\partial g+\rho(g)\partial f$.  The associated
Hodge operator is
\begin{equation}
 D_{m,c}=\begin{pmatrix}0&\partial^*\\\partial&0\end{pmatrix}.
 \label{eq:supp-graph-Dirac}
\end{equation}

\begin{theorem}[Renewal Hodge square and exact Lipschitz seminorm]
\label{thm:supp-graph-square}
Let $\mathcal L^\dagger$ be the adjoint of $\mathcal L$ in $\ell^2(V,m)$ and
let $\mathcal L_s=(\mathcal L+\mathcal L^\dagger)/2$.  Then
\begin{equation}
 {D_{m,c}^2|_{\Hh_0}=\partial^*\partial=-2\mathcal L_s.}
 \label{eq:supp-graph-square}
\end{equation}
For every real $f\in C(V)$,
\begin{equation}
 {
 L_{m,c}(f)^2
 =\norm{[D_{m,c},f]}^2
 =\max_{x\in V}\frac1{m_x}
   \sum_yc_{xy}|f(y)-f(x)|^2.}
 \label{eq:supp-graph-Lip}
\end{equation}
The kernel of $L_{m,c}$ is the constants, so the Connes formula gives a finite
metric on $V$.  On the reversible branch, $\mathcal L=\mathcal L_s$ and
$D_{m,c}^2|_{\Hh_0}=-2\mathcal L$.
\end{theorem}

\begin{proof}
The adjoint generator has rates $q_{yx}/m_x$, so $\mathcal L_s$ has rates
$c_{xy}/m_x$.  Therefore
\[
 -\ip{f}{\mathcal L_sf}_0
 =\frac12\sum_{x,y}c_{xy}|f(y)-f(x)|^2
 =\frac12\norm{\partial f}_1^2.
\]
Polarization gives \eqref{eq:supp-graph-square}.  The lower-left commutator
block satisfies
\[
 (B_fg)(x,y)=(f(y)-f(x))g(x).
\]
Hence
\[
 \norm{B_fg}_1^2
 =\sum_x|g(x)|^2\sum_yc_{xy}|f(y)-f(x)|^2.
\]
Its norm relative to $\sum_xm_x|g(x)|^2$ is the maximum in
\eqref{eq:supp-graph-Lip}.  Connectedness makes $\partial f=0$ equivalent to
constancy.
\end{proof}

The connectedness assumption has a state-space consequence as well as a
pointwise one.  Proposition~\ref{prop:finite-CQMS} implies that
$(C(V)_{\rm sa},L_{m,c})$ is a compact quantum metric space.  This
finite-dimensional Lip-norm fact is standard and is not the new assertion of this
section.  The relevant structural identity is the exact equality
\eqref{eq:supp-graph-Lip} between the commutator seminorm and the local
carr\'e-du-champ of the symmetric Markov generator; the continuum result below is
the uniform convergence of the resulting $A_3$ Connes metric to the independently
reconstructed Riemannian metric.  Its states are
probability distributions on $V$, and the spectral metric on those distributions
induces their usual weak-$*$ topology.  This conclusion does not require the
Hodge operator itself to be invertible; it is the scalar kernel of the
\emph{commutator seminorm} that matters.  The following large one-form kernel is
therefore an issue for the operator continuum limit, not for the finite Lip-norm.

\begin{proposition}[Extensive kernel of the one-step graph calculus]
\label{cth:supp-extensive-kernel}
Let the connected underlying graph have $n$ vertices and $e$ undirected edges,
and retain both orientations in $\Hh_1$.  Then
\begin{equation}
 {\dim\Ker D_{m,c}=2e-n+2.}
 \label{eq:supp-extensive-kernel}
\end{equation}
On a periodic $A_3$ root graph, $e=6n$, so the kernel has dimension $11n+2$.
Thus the one-step universal graph calculus cannot itself have a compact-
resolvent cofinal limit retaining all one-form directions.
\end{proposition}

\begin{proof}
Connectedness gives $\dim\Ker\partial=1$ and $\rank\partial=n-1$.  Since the
directed one-form space has dimension $2e$,
$\dim\Ker\partial^*=2e-n+1$.  Because
$\Ker D=(\Ker\partial)\oplus(\Ker\partial^*)$, the formula follows.
\end{proof}

\subsection{The periodic $A_3$ model and convergence of the Connes metric}

Let
\begin{equation}
 W_0=\left\{x\in\R^4:\sum_{i=1}^4x_i=0\right\},
 \qquad
 \Lambda_{A_3}=W_0\cap\Z^4,
 \label{eq:supp-A3-lattice}
\end{equation}
and let
$\Phi_{A_3}=\{e_i-e_j:i\ne j\}$ be the twelve oriented roots.

\begin{lemma}[$A_3$ tight-frame identity]
\label{lem:supp-A3-frame}
On $W_0$,
\begin{equation}
 {\sum_{\alpha\in\Phi_{A_3}}\alpha\alpha^{\mathsf T}=8I_{W_0}.}
 \label{eq:supp-A3-frame}
\end{equation}
Equivalently,
$\sum_{i\ne j}(v_i-v_j)^2=8\norm v^2$ for $v\in W_0$.
\end{lemma}

\begin{proof}
Expand the left side.  Each $v_i^2$ occurs six times in the two diagonal sums,
while
$\sum_{i\ne j}v_iv_j=(\sum_iv_i)^2-\sum_iv_i^2=-\sum_iv_i^2$.
The total is $8\sum_i v_i^2$.  Polarization gives the operator identity.
\end{proof}

Put
\begin{equation}
 M=W_0/\Lambda_{A_3},
 \qquad
 V_h=h\Lambda_{A_3}/\Lambda_{A_3},
 \qquad h=d^{-1},
 \label{eq:supp-A3-torus}
\end{equation}
and assign each root edge the reversible rate and vertex mass
\begin{equation}
 k_{\alpha,h}=\frac1{8h^2},
 \qquad
 m_h(x)=h^3,
 \qquad
 c_{\alpha,h}=\frac h8.
 \label{eq:supp-A3-data}
\end{equation}

\begin{theorem}[Exact finite $A_3$ spectralization]
\label{thm:supp-A3-finite}
The graph spectral triple has
\begin{equation}
 {
 L_h(f)^2
 =\max_{x\in V_h}\frac1{8h^2}
  \sum_{\alpha\in\Phi_{A_3}}
  |f(x+h\alpha)-f(x)|^2,}
 \label{eq:supp-A3-Lip}
\end{equation}
\begin{equation}
 {D_h^2|_{C(V_h)}=-2\mathcal L_h,}
 \qquad
 (\mathcal L_hf)(x)
 =\frac1{8h^2}\sum_{\alpha\in\Phi_{A_3}}
   (f(x+h\alpha)-f(x)).
 \label{eq:supp-A3-square}
\end{equation}
For every $F\in C^2(M)$,
\begin{equation}
 \left|L_h(F|_{V_h})-\norm{\nabla F}_{L^\infty(M)}\right|
 \le C_{A_3}h\norm F_{C^2(M)}.
 \label{eq:supp-A3-Lip-convergence}
\end{equation}
Moreover $\mathcal L_hF\to\frac12\Delta F$ uniformly for $F\in C^4(M)$.
\end{theorem}

\begin{proof}
The first two identities follow from \Cref{thm:supp-graph-square} and
\eqref{eq:supp-A3-data}.  Taylor expansion gives
\[
 \frac{F(x+h\alpha)-F(x)}h
 =\alpha\cdot\nabla F(x)+r_{\alpha,h}(x),
 \qquad |r_{\alpha,h}(x)|\le Ch\norm F_{C^2}.
\]
Apply the reverse triangle inequality in $\ell^2(\Phi_{A_3})$ and
\Cref{lem:supp-A3-frame}, then maximize over the mesh, to obtain
\eqref{eq:supp-A3-Lip-convergence}.  For the generator, the linear terms cancel
by central symmetry.  The quadratic term is
$\frac1{16}\sum_\alpha(\alpha\otimes\alpha):\nabla^2F=\frac12\Delta F$;
the cubic terms cancel and the remainder is $O(h^2\norm F_{C^4})$.
\end{proof}

\begin{lemma}[Compactness of discrete unit-Lipschitz functions]
\label{lem:supp-A3-compactness}
Let $h_n\downarrow0$ and let $f_n:V_{h_n}\to\R$ satisfy
$L_{h_n}(f_n)\le1$ and $f_n(x_{0,n})=0$ for $x_{0,n}\to x_0$.  Then a
subsequence of periodic piecewise-affine interpolants converges uniformly to
$F\in W^{1,\infty}(M)$ with
\begin{equation}
 |\nabla F|\le1
 \quad\text{almost everywhere}.
 \label{eq:supp-A3-limit-Lip}
\end{equation}
\end{lemma}

\begin{proof}
Equation~\eqref{eq:supp-A3-Lip} gives the root-edge bound
$|f_n(x+h_n\alpha)-f_n(x)|\le\sqrt8h_n$.  A fixed shape-regular periodic
simplicial refinement of the $A_3$ lattice therefore gives equi-Lipschitz,
uniformly bounded interpolants.  Arzel\`a--Ascoli gives uniform convergence
along a subsequence.

The difference quotients
$g_{n,\alpha}=(f_n(\cdot+h_n\alpha)-f_n)/h_n$ lie pointwise in the closed convex
ball $\frac18\sum_\alpha|z_\alpha|^2\le1$.  Passing weak-* in $L^\infty$ and
using discrete periodic summation by parts gives
$g_\alpha=\alpha\cdot\nabla F$ distributionally.  The tight-frame identity then
yields
$|\nabla F|^2=\frac18\sum_\alpha|\alpha\cdot\nabla F|^2\le1$.
\end{proof}

\begin{theorem}[Renewal--spectral distance identity]
\label{thm:supp-A3-distance}
Let $d_h$ be the pure-state Connes distance defined by
\eqref{eq:supp-A3-Lip}, and let $d_g$ be the flat geodesic distance of the
metric independently reconstructed from the renewal second moment.  Then
\begin{equation}
 {
 \sup_{x,y\in V_h}|d_h(x,y)-d_g(x,y)|\longrightarrow0.}
 \label{eq:supp-A3-distance}
\end{equation}
\end{theorem}

\begin{proof}
For the lower bound, mollify the one-Lipschitz function
$z\mapsto d_g(z,y)$ periodically to obtain smooth $F_{y,\varepsilon}$ with
$\|F_{y,\varepsilon}-d_g(\cdot,y)\|_\infty\le\varepsilon$ and
$\|\nabla F_{y,\varepsilon}\|_\infty\le1$, uniformly in $y$.
By \eqref{eq:supp-A3-Lip-convergence}, its sampled seminorm is at most
$1+o_h(1)$.  The dual formula then gives
$\liminf_h(d_h-d_g)\ge-2\varepsilon$ uniformly; let $\varepsilon\downarrow0$.

If the upper bound failed, there would be $h_n\downarrow0$, vertices
$x_n,y_n$, and real $f_n$ with $L_{h_n}(f_n)\le1$ and
$f_n(x_n)-f_n(y_n)\ge d_g(x_n,y_n)+\varepsilon$.  Normalize by
$f_n(y_n)=0$ and pass to subsequences.  By
\Cref{lem:supp-A3-compactness}, the interpolants converge uniformly to a
one-Lipschitz $F$.  Thus
$F(x)-F(y)\le d_g(x,y)$, contradicting the limiting inequality.
\end{proof}


Each finite $A_3$ approximant is consequently a compact quantum metric space.
The distance convergence just proved concerns its pure states, that is, the
vertices.  It is not by itself a quantum Gromov--Hausdorff, propinquity, or
spectral-propinquity convergence theorem: those stronger notions compare the full
Lip-norms, state spaces, and, in the spectral setting, the corresponding operator
data \cite{Rieffel2004,Latremoliere2022,FarsiLatremolierePacker2024}.  In
particular, a quadratic graph commutator seminorm need not be the ordinary
Lipschitz seminorm of its pure-state metric \cite{Rieffel1999}.  Nor does the finite-dimensional conclusion of
Proposition~\ref{prop:finite-CQMS} automatically survive an infinite-dimensional
limit.  A limiting quantum compact metric-space statement would require a
uniform total-boundedness or equicontinuity estimate for normalized Lipschitz
balls (the noncommutative analogue of the Arzel\`a--Ascoli input), together with
control of the limiting seminorm.  Lemma~\ref{lem:supp-A3-compactness} supplies
the required precompactness for the scalar functions used in the pure-state
metric proof, but it is not asserted here to prove convergence of the full
quantum state spaces.  Keeping these assertions separate avoids conflating a
pointwise metric limit with a noncommutative metric-space limit.

\subsection{Local spin discretization and norm-resolvent convergence under atlas stability}

The periodic model makes the metric comparison explicit.  We next formulate the
local construction on an arbitrary closed Riemannian spin manifold
$(M^d,g,\mathfrak s)$ and separate what is proved chartwise from the additional
global compatibility needed for norm-resolvent convergence.  The local scalar
ingredient is a positive decomposition of the inverse metric.

\begin{lemma}[Positive local stencil]
\label{lem:positive-stencil-main}
Let
\[
 \mathcal R_d^+
 =\{e_i:1\le i\le d\}
 \cup\{e_i+e_j,e_i-e_j:1\le i<j\le d\}.
\]
After a linear normalization at a point, every inverse metric sufficiently close
to the identity admits a decomposition
\begin{equation}
 g^{-1}(x)=\sum_{r\in\mathcal R_d^+}w_r(x)rr^{\mathsf T}
 \label{eq:positive-stencil-main}
\end{equation}
with all $w_r(x)>0$.  Hence a finite good cover of a closed Riemannian manifold may
be chosen so that such a positive stencil exists on every chart.
\end{lemma}

\begin{proof}
Write $A=(A_{ij})$ for the normalized inverse metric and fix $\eta>0$ such that
$|A_{ij}|<\eta$ for $i\neq j$ and $A_{ii}>(d-1)\eta$.  Set
\[
 w_i=A_{ii}-(d-1)\eta,
 \qquad
 w_{ij}^{\pm}=\tfrac12(\eta\pm A_{ij}).
\]
All coefficients are positive and direct expansion gives
\[
 A=\sum_iw_ie_ie_i^{\mathsf T}
 +\sum_{i<j}\bigl[
 w_{ij}^+(e_i+e_j)(e_i+e_j)^{\mathsf T}
 +w_{ij}^-(e_i-e_j)(e_i-e_j)^{\mathsf T}\bigr].
\]
Continuity and compactness then give a finite cover.
\end{proof}

The weights in \eqref{eq:positive-stencil-main} define reversible local jump
processes whose second moments converge to $g^{-1}$ and whose generators converge
to $\frac12\Delta_g$.  To recover spin geometry one chooses the marked oriented
orthonormal coframe and its spin lift.  This is additional information: the scalar
metric bracket alone does not determine the spin structure.  With
$c^j=E_a{}^j\gamma_a$ and $\rho=\sqrt{\det g}$, the geometric Dirac operator admits
the density-symmetric representation
\begin{equation}
 \rho^{1/2}D_g\rho^{-1/2}
 =-\frac{i}{2}\sum_{j=1}^d
 \bigl(M_{c^j}\nabla_j+\nabla_jM_{c^j}\bigr),
 \label{eq:density-symmetric-main}
\end{equation}
which is well suited to self-adjoint finite-difference discretization
\cite{lawson-michelsohn}.

A naive centered lattice Dirac operator generally has extra low-energy species;
we therefore use a covariant Wilson term, as in the standard resolution of the
fermion-doubling problem \cite{nielsenninomiya1981a,nielsenninomiya1981b}.
Its role in rigorous continuum limits is discussed in \cite{Nakamura2024}.
For fixed positive Wilson parameter, this term is $O(h)$ on sampled smooth
spinors but of size $O(h^{-1})$ at nonzero lattice corners.  It removes the
unwanted ultraviolet species without adding a term to the limiting equation
on the smooth core.  The distinction between these two estimates is essential:
the Wilson term does not tend to zero in operator norm on the whole lattice
space.  A fixed two-dimensional normal Clifford factor supplies the stable
real-even convention.

Passing from local coordinate estimates to a closed manifold requires more
than a positive overlap Gram.  The assembled finite operators must have a
uniform first-order graph estimate, compatible reconstructions into the fixed
spin bundle, and consistency on its smooth core.  These requirements are
stated precisely in Definition~\ref{def:stable-spin-atlas}.  They also specify
which self-adjoint continuum realization is intended: the closure on smooth
spinors, with domain $H^1$, rather than independent boundary realizations on
the coordinate charts.

\begin{theorem}[Norm-resolvent convergence criterion under atlas stability]
\label{thm:compact-spin-main}
Let $(M^d,g,\mathfrak s)$ be a closed Riemannian spin manifold, and suppose that
its finite covariant Wilson operators form a compatible stable spin
discretization in the sense of Definition~\ref{def:stable-spin-atlas}.
Then, for every $z\in\C\setminus\R$,
\begin{equation}
 \left\|W_h(D_h-z)^{-1}W_h^*
       -(\widehat D_{M,\mathfrak s}-z)^{-1}\right\|\longrightarrow0,
 \qquad \widehat D_{M,\mathfrak s}=D_{M,\mathfrak s}\otimes\sigma_1.
 \label{eq:compact-spin-resolvent-main}
\end{equation}
The real and grading structures are retained when intertwined as specified
there, and the global construction uses the fixed spin structure and
orientation.  Independently, every closed Riemannian manifold admits finite
metric triples on refining meshes with
\begin{equation}
 \sup_{x,y\in V_h}|d_h^{\rm met}(x,y)-d_g(x,y)|\longrightarrow0.
 \label{eq:compact-spin-distance-main}
\end{equation}
Thus the spin conclusion is a convergence theorem for compatible stable
assemblies; the scalar metric approximation does not require those additional
spin-assembly hypotheses.
\end{theorem}

\begin{proof}
Lemma~\ref{lem:positive-stencil-main} supplies reversible local generators
whose scalar brackets approximate $g^{-1}$.  On the spin side, the frozen
Wilson square controls the full forward-difference symbol uniformly, including
nonzero Brillouin corners.  Appendix~\ref{app:compact-spin} gives the explicit
symbol comparison and derives the chartwise G\aa rding bound by localization
on fixed, sufficiently small coordinate patches.

In a compatible global assembly, smooth cutoffs with
$\sum_\alpha\chi_\alpha^2=1$ introduce no order-one boundary term:
$\sum_\alpha\chi_\alpha[\widehat D,\chi_\alpha]=0$ after the spin and
density transitions are included.  The analogous Wilson double-commutator
error is $O(h)$ on common meshes.  For independently meshed charts, the
interpolation and assembled stability requirements are part of the explicit
atlas hypotheses, not a consequence of positivity alone.

Self-adjointness bounds the discrete resolvents by
$|\operatorname{Im}z|^{-1}$.  The global graph estimate and reconstruction
bounds make their images collectively compact in $L^2$.  Testing the
resolvent equation against sampled smooth spinors identifies every limit with
$(\widehat D-z)^{-1}$; the same argument at $\bar z$ gives adjoint
convergence.  These facts yield norm convergence by the compactness argument
proved in Theorem~\ref{thm:supp-compact-spin-resolvent}.

For the metric conclusion, connect mesh points at distances of a mesoscopic
scale $\rho_h\downarrow0$, with $h/\rho_h\to0$, using endpoint blocks weighted
by the lengths of short Riemannian curves.  The commutator seminorm is the
ordinary graph Lipschitz seminorm.  Approximating minimizing geodesics by mesh
paths gives, uniformly,
\[
 d_g(x,y)\le d_h^{\rm met}(x,y)
 \le d_g(x,y)+C\left(\rho_h+\frac h{\rho_h}\right).
\]
This proves \eqref{eq:compact-spin-distance-main}; the complete metric
construction is in Theorem~\ref{thm:supp-global-metric}.
\end{proof}

\begin{remark}
The theorem is Riemannian.  It does not require, and does not assert, a Lorentzian
or interacting quantum-field-theoretic completion.  The role of the fixed normal
Clifford factor is only to provide a stable real-even lattice realization.
The local construction proves the necessary chartwise estimates.  An
unconditional assertion that a particular atlas construction realizes every
closed spin manifold additionally requires verification of the global
compatibility and stability properties in Definition~\ref{def:stable-spin-atlas}.
The convergence criterion does not assume that this follows merely from an
overlap Gram quotient.
\end{remark}


% =============================================================================
\section{Nonfaithfulness and inverse fibres}
\label{sec:inverse}
% =============================================================================

The preceding sections show that the spectralization map has a broad essential
image.  It is nevertheless far from injective.  Everything in this section is
formulated in standard operator-theoretic terms once the graph generator, the
operator-valued self-energy, or the marked spectral packet is specified; the
Renewal terminology records one motivating interpretation and is not an
additional hypothesis.  We describe three independent mechanisms of
nonfaithfulness: directed current at fixed symmetric metric data, hidden
dynamical memory, and global marked monodromy.  Each admits an intrinsic
coordinate, and the principal coordinates transform functorially under coarse
graining.

\subsection{Directed currents at fixed metric geometry}

Fix a connected undirected graph $G=(V,E)$, positive masses $m_x$, and positive
conductances $c_e$.  Choose one orientation of each undirected edge and let
$B:\R^E\to\R^V$ be the incidence matrix.  An oriented current is extended to
the reverse orientation by $j_{\bar e}=-j_e$.

\begin{definition}[Admissible cycle-current polytope]
\label{def:supp-current-polytope}
Define
\begin{equation}
 \mathcal J(G,c)
 =\{j\in\Ker B:\abs{j_e}<c_e\text{ for every }e\in E\}.
 \label{eq:supp-current-polytope}
\end{equation}
Allowing non-strict inequalities gives the closed fibre in which some directed
rates may vanish.
\end{definition}

\begin{theorem}[Stationary Markov realization fibre of a metric graph triple]
\label{thm:supp-graph-realization}
Let
$\mathfrak T(G,m,c)=(C(V),\Hh_{m,c},D_{m,c},J,\gamma)$ be the metric graph
spectral triple of \Cref{thm:supp-graph-square}.  Then:
\begin{enumerate}[label=\textup{(R\arabic*)},leftmargin=3em]
\item it has the reversible Markov realization
      \begin{equation}
      k_{xy}^{\rm rev}=\frac{c_{xy}}{m_x};
      \label{eq:supp-reversible-realization}
      \end{equation}
\item every $j\in\mathcal J(G,c)$ gives a stationary Markov realization
      \begin{equation}
      q_{xy}=c_{xy}+j_{xy},
      \qquad
      k_{xy}=\frac{c_{xy}+j_{xy}}{m_x};
      \label{eq:supp-current-realization}
      \end{equation}
\item conversely, every stationary renewal generator with masses $m$ and
      symmetric conductances $c$ is uniquely of the form
      \eqref{eq:supp-current-realization};
\item every realization in the fibre has the same metric spectral triple.
\end{enumerate}
Thus, up to unread parallel-channel refinements,
\begin{equation}
 {
 \RReal_{\rm graph}(\mathfrak T(G,m,c))
 \cong\mathcal J(G,c).}
 \label{eq:supp-realization-fibre}
\end{equation}
\end{theorem}

\begin{proof}
For \eqref{eq:supp-reversible-realization}, the stationary flux is
$q_{xy}=c_{xy}=q_{yx}$, so detailed balance holds.  For
\eqref{eq:supp-current-realization}, positivity follows from
$|j_{xy}|<c_{xy}$.  Moreover
\[
 \sum_y(q_{xy}-q_{yx})=2\sum_yj_{xy}=0
\]
because $Bj=0$, so $m$ is stationary.  Conversely, for a stationary continuous-time Markov
generator define $j=(q-q^{\mathsf T})/2$.  Stationarity gives $Bj=0$, positivity gives
the edge inequalities, and $q=c+j$ is unique.  The graph spectralization uses
only $m$ and $(q+q^{\mathsf T})/2=c$.
\end{proof}

\begin{corollary}[Dimension and faithfulness of the fibre]
\label{cor:supp-fibre-dimension}
For connected $G$,
\begin{equation}
 {\dim\mathcal J(G,c)=|E|-|V|+1=b_1(G).}
 \label{eq:supp-fibre-dimension}
\end{equation}
The reversible realization is unique exactly when $G$ is a tree.  If
$b_1(G)>0$, the metric spectralization map is not faithful.
\end{corollary}

\begin{proof}
The incidence matrix of a connected graph has rank $|V|-1$, so
$\dim\Ker B=|E|-|V|+1$.  The strict inequalities cut out an open neighborhood
of zero in this cycle space.  A connected graph has zero first Betti number
exactly when it is a tree.
\end{proof}

\begin{proposition}[Entropy production inside one spectral fibre]
\label{prop:supp-entropy}
For a stationary realization with current $j$, the steady entropy production
is
\begin{equation}
 {
 \sigma(j)
 =2\sum_{e\in E}j_e
   \log\frac{c_e+j_e}{c_e-j_e}\ge0.}
 \label{eq:supp-entropy}
\end{equation}
It vanishes exactly at $j=0$.
\end{proposition}

\begin{proof}
Start from
$\frac12\sum_{x,y}(q_{xy}-q_{yx})\log(q_{xy}/q_{yx})$, pair the two
orientations of each edge, and substitute $q_e=c_e+j_e$,
$q_{\bar e}=c_e-j_e$.  Each summand is nonnegative because $j_e$ and the
logarithm have the same sign; it vanishes only for $j_e=0$.
\end{proof}

\begin{example}[Three-cycle universality family]
On an oriented three-cycle with equal conductance $c$, every constant current
$j_e=s$, $|s|<c$, is divergence free.  All values of $s$ give one metric
spectral triple, whereas
\[
 \sigma(s)=6s\log\frac{c+s}{c-s}
\]
varies continuously from zero and diverges as $|s|\uparrow c$.
\end{example}

\begin{theorem}[Finite spectral universality]
\label{thm:supp-spectral-universality}
There exist nonisomorphic stationary predictive graph dynamics
$T_1\not\cong T_2$ with isomorphic metric spectral images.  More strongly, for
every connected graph with $b_1(G)>0$, the fibre contains a
$b_1(G)$-dimensional open family of dynamically inequivalent renewal laws.
\end{theorem}

\begin{proof}
Choose two distinct currents in \eqref{eq:supp-current-polytope}.  Their
directed fluxes and path probabilities differ, and generically their entropy
productions differ, so the history-preserving renewal packets are not
isomorphic.  \Cref{thm:supp-graph-realization} gives identical metric spectral
images, and \Cref{cor:supp-fibre-dimension} gives the dimension.
\end{proof}

\subsection{The modular coordinate of the current fibre}

For $j\in\mathcal J(G,c)$ define the real affinity
\begin{equation}
 a_e=\operatorname{arctanh}\frac{j_e}{c_e}
 =\frac12\log\frac{c_e+j_e}{c_e-j_e}.
 \label{eq:supp-affinity}
\end{equation}
Exact one-cochains are $\Ran B^{\mathsf T}$, so $a$ has a class
$[a]\in H^1(G;\R)=\R^E/\Ran B^{\mathsf T}$.

\begin{lemma}[Zero affinity class forces reversibility]
\label{lem:supp-zero-affinity}
If $j\in\mathcal J(G,c)$ is divergence free and $[a]=0$, then $j=0$.
\end{lemma}

\begin{proof}
Write $a=B^{\mathsf T}\phi$.  Since
$j_e=c_e\tanh((B^{\mathsf T}\phi)_e)$ and $Bj=0$,
\[
 0=\langle\phi,Bj\rangle
 =\sum_ec_e(B^{\mathsf T}\phi)_e
       \tanh((B^{\mathsf T}\phi)_e).
\]
Every summand is nonnegative and vanishes only when the edge difference is
zero.  Connectedness makes $\phi$ constant, hence $a=j=0$.
\end{proof}

\begin{theorem}[Nonlinear Hodge theorem for stationary affinities]
\label{thm:supp-nonlinear-Hodge}
The map
\begin{equation}
 \mathfrak A:\mathcal J(G,c)\longrightarrow H^1(G;\R),
 \qquad
 j\longmapsto[\operatorname{arctanh}(j/c)]
 \label{eq:supp-affinity-map}
\end{equation}
is a diffeomorphism.  More explicitly, for a class $\alpha$ and representative
$a_0$, define on mean-zero vertex potentials
\begin{equation}
 \mathcal F_\alpha(\phi)
 =\sum_{e\in E}c_e
  \log\cosh(a_{0,e}+(B^{\mathsf T}\phi)_e).
 \label{eq:supp-convex-affinity-functional}
\end{equation}
This functional is coercive and strictly convex.  Its unique minimizer
$\phi_\alpha$ gives
\begin{equation}
 a_\alpha=a_0+B^{\mathsf T}\phi_\alpha,
 \qquad
 j_{\alpha,e}=c_e\tanh(a_{\alpha,e}),
 \qquad Bj_\alpha=0,
 \label{eq:supp-affinity-inverse}
\end{equation}
and this is the unique current with modular class $\alpha$.
\end{theorem}

\begin{proof}
On mean-zero potentials, connectedness gives a Poincar\'e estimate
$\|\phi\|\le C\|B^{\mathsf T}\phi\|$.  Since
$\log\cosh t\ge|t|-\log2$, the functional is coercive.  Its Hessian is
\[
 B\,\diag\!\left(c_e\operatorname{sech}^2
  (a_{0,e}+(B^{\mathsf T}\phi)_e)\right)B^{\mathsf T},
\]
which is positive definite on the mean-zero subspace.  Thus there is a unique
minimizer.  Its Euler equation is
\[
 0=\sum_ec_e\tanh(a_{\alpha,e})(B^{\mathsf T}\psi)_e
  =\langle Bj_\alpha,\psi\rangle,
\]
so $j_\alpha$ is divergence free.  Conversely, if $j$ is stationary and
$a=\operatorname{arctanh}(j/c)$, then zero is the critical point of
\eqref{eq:supp-convex-affinity-functional} in the class $[a]$; strict
convexity makes it the unique minimizer.  Smooth dependence follows from the
implicit-function theorem and the positive Hessian.
\end{proof}

\begin{corollary}[Faithful metric--modular enrichment]
\label{cor:supp-faithful-modular}
For fixed $(G,m,c)$, the enriched object
\begin{equation}
 \mathfrak S_{\rm mod}(T)
 =\bigl(C(V),\Hh_{m,c},D_{m,c},J,\gamma;[a]\bigr)
 \label{eq:supp-modular-enrichment}
\end{equation}
determines the complete stationary generator through
\[
 [a]\longmapsto a_\alpha
 \longmapsto j=c\tanh a_\alpha
 \longmapsto q=c+j
 \longmapsto k_{xy}=q_{xy}/m_x.
\]
Hence the modular enrichment is faithful on each finite graph fibre.
\end{corollary}

\begin{proof}
Apply \Cref{thm:supp-nonlinear-Hodge} and
\Cref{thm:supp-graph-realization}.
\end{proof}

The entropy production is the modular pairing
\begin{equation}
 {\sigma(j)=4\sum_{e\in E}j_ea_e,}
 \label{eq:supp-entropy-modular-pairing}
\end{equation}
and is unchanged by $a\mapsto a+B^{\mathsf T}\phi$ because
$\langle j,B^{\mathsf T}\phi\rangle=\langle Bj,\phi\rangle=0$.

\subsection{Relative modular data and cellular curvature}

Let $E^{\rm or}$ contain both orientations of every edge, with reversal
$e\mapsto\bar e$, source $s(e)$, terminal $t(e)=s(\bar e)$, and positive
directed stationary flux $q_e$.  On $\Hh_E=\ell^2(E^{\rm or})$, define
\begin{equation}
 (J_E\xi)(e)=\overline{\xi(\bar e)},
 \qquad
 (\Delta_q\xi)(e)=\frac{q_e}{q_{\bar e}}\xi(e)=e^{2a_e}\xi(e),
 \qquad
 S_q=J_E\Delta_q^{1/2}.
 \label{eq:supp-Tomita-packet}
\end{equation}
Source and terminal multiplication are
$(\lambda_E(f)\xi)(e)=f(s(e))\xi(e)$ and
$(\rho_E(f)\xi)(e)=f(t(e))\xi(e)$.

\begin{theorem}[Finite modular-reversal identities]
\label{thm:supp-Tomita}
The packet \eqref{eq:supp-Tomita-packet} satisfies
\begin{equation}
 {
 J_E^2=I,
 \qquad
 J_E\Delta_qJ_E=\Delta_q^{-1},
 \qquad
 S_q^2=I,}
 \label{eq:supp-Tomita-identities}
\end{equation}
\begin{equation}
 {
 J_E\lambda_E(f^*)J_E=\rho_E(f),
 \qquad
 S_q\lambda_E(f^*)S_q^{-1}=\rho_E(f),
 \qquad
 \log\Delta_q(e)=2a_e.}
 \label{eq:supp-Tomita-opposite}
\end{equation}
The reversible branch is exactly $\Delta_q=I$.
\end{theorem}

\begin{proof}
Reversal preserves counting measure, so $J_E$ is antiunitary and squares to
the identity.  Since $a_{\bar e}=-a_e$, conjugating the diagonal multiplier
$e^{2a_e}$ by $J_E$ gives its inverse; hence $S_q^2=I$.  The source of the
reversed edge is the terminal of the original edge, proving the opposite-action
identities.  The logarithm statement is the definition of $a_e$.
\end{proof}

For an oriented cycle $C$,
\begin{equation}
 \operatorname{Hol}_{\Delta}(C)
 =\prod_{e\in C}\left(\frac{q_e}{q_{\bar e}}\right)^{\epsilon_{C,e}}
 =\exp\left(2\sum_{e\in C}\epsilon_{C,e}a_e\right).
 \label{eq:supp-modular-cycle-holonomy}
\end{equation}

Now let
\begin{equation}
 C^0\xrightarrow{d_0}C^1\xrightarrow{d_1}C^2
 \xrightarrow{d_2}C^3
 \label{eq:supp-cellular-Hodge-complex}
\end{equation}
be a finite real Hilbert complex and let
$\mathscr H^1=\Ker d_1\cap\Ker d_0^*$.

\begin{theorem}[Cellular modular Hodge reconstruction]
\label{thm:supp-modular-Hodge}
For $a\in C^1$, define
\begin{equation}
 h_a=P_{\mathscr H^1}a,
 \qquad
 F_a=d_1a.
 \label{eq:supp-modular-hF}
\end{equation}
Then $a$ has the unique orthogonal decomposition
\begin{equation}
 {a=d_0\phi+h_a+d_1^*\psi.}
 \label{eq:supp-modular-Hodge}
\end{equation}
Its coexact part is reconstructed from the curvature by
\begin{equation}
 P_{\Ran d_1^*}a
 =d_1^*(d_1d_1^*)^\dagger F_a.
 \label{eq:supp-modular-coexact}
\end{equation}
Consequently
\begin{equation}
 C^1/\Ran d_0\longrightarrow\mathscr H^1\oplus\Ran d_1,
 \qquad
 [a]_{\rm gauge}\longmapsto(h_a,F_a)
 \label{eq:supp-modular-gauge-map}
\end{equation}
is injective and onto the displayed product.  A canonical representative is
\begin{equation}
 a_{h,F}=h+d_1^*(d_1d_1^*)^\dagger F.
 \label{eq:supp-modular-representative}
\end{equation}
Moreover $d_2F_a=0$ and
\begin{equation}
 \norm a^2
 =\norm{d_0\phi}^2+\norm{h_a}^2
 +\ip{F_a}{(d_1d_1^*)^\dagger F_a}.
 \label{eq:supp-modular-norm}
\end{equation}
\end{theorem}

\begin{proof}
Finite Hodge decomposition gives
$C^1=\Ran d_0\oplus\mathscr H^1\oplus\Ran d_1^*$, proving
\eqref{eq:supp-modular-Hodge}.  Since $d_1$ annihilates the first two
summands, $F_a=d_1d_1^*\psi$; Moore--Penrose inversion on $\Ran d_1$ gives
\eqref{eq:supp-modular-coexact}.  If two cochains have the same $(h,F)$, their
difference has zero harmonic projection and lies in $\Ker d_1$.  Its coexact
part then has zero norm, so the difference is exact.  This proves injectivity;
\eqref{eq:supp-modular-representative} proves surjectivity.  Orthogonality
gives the norm identity, and $d_2d_1=0$ gives the Bianchi identity.
\end{proof}

\begin{corollary}[Flat, globally twisted, and curved modular branches]
\label{cor:supp-modular-trichotomy}
For a cellular affinity $a$:
\begin{enumerate}[label=\textup{(M\arabic*)},leftmargin=3em]
\item if $F_a\ne0$, the modular transport has local face curvature;
\item if $F_a=0$ but $h_a\ne0$, it is locally flat with nontrivial global positive
      holonomy;
\item if $F_a=h_a=0$, then $a=d_0\phi$ and the modular packet is removed by a
      positive vertex gauge.
\end{enumerate}
On the periodic three-torus, the global flat sector has dimension three.
\end{corollary}

\begin{theorem}[Face holonomy]
\label{thm:supp-face-holonomy}
For every oriented loaded two-cell $p$,
\begin{equation}
 {
 \prod_{e\subset\partial p}
 \left(\frac{q_e}{q_{\bar e}}\right)^{\epsilon_{p,e}}
 =\exp(2F_a(p)).}
 \label{eq:supp-face-holonomy}
\end{equation}
Thus a source--terminal modular transport descends to a flat cell-line packet
exactly when every face-curvature residual vanishes.
\end{theorem}

\begin{proof}
By the cellular coboundary formula,
$F_a(p)=\sum_e\epsilon_{p,e}a_e$.  Substitute
$q_e/q_{\bar e}=e^{2a_e}$ and multiply.
\end{proof}
\subsection{Hidden dynamical memory}
\label{subsec:dynamic-memory}

The graph-current fibre concerns directed dynamics at fixed metric geometry.
A second inverse fibre appears when auxiliary memory degrees of freedom have been eliminated
into an effective head operator.  The dynamic coordinate is the whole
resolvent self-energy rather than one static Schur value.

\begin{definition}[Positive operator-valued tail measure]
\label{def:supp-POVM}
Let $H$ be finite dimensional.  A bounded positive operator-valued tail
measure is a weakly countably additive map
\[
 \mu:\mathfrak B((0,\infty))\to\mathcal B(H),
 \qquad
 \mu(E)\succeq0,
 \qquad
 \|\mu((0,\infty))\|<\infty.
\]
Its Stieltjes self-energy is
\begin{equation}
 {
 \Sigma_\mu(z)
 =\int_{(0,\infty)}
  \frac{1}{\lambda-z}\,\dd\mu(\lambda),
 \qquad z\in\mathbb C\setminus[0,\infty).}
 \label{eq:supp-Stieltjes}
\end{equation}
\end{definition}

\begin{construction}[Canonical minimal Stieltjes dilation]
\label{con:supp-Naimark-tail}
Let $\mathscr S_\mu$ be the finite $H$-valued Borel simple functions and set
\[
 \left\langle
 \sum_j\one_{E_j}\xi_j,
 \sum_k\one_{F_k}\eta_k
 \right\rangle_\mu
 =
 \sum_{j,k}
 \langle\xi_j,\mu(E_j\cap F_k)\eta_k\rangle.
\]
Quotient by the null space and complete to $K_\mu$.  Multiplication by
$\one_E$ defines a projection-valued measure $P_\mu(E)$, and
\[
 C_\mu=\int\lambda\,\dd P_\mu(\lambda),
 \qquad
 V_\mu\xi=[\one_{(0,\infty)}\xi],
 \qquad
 B_\mu=V_\mu^*.
\]
\end{construction}

\begin{theorem}[Operator-measure realization and minimality]
\label{thm:supp-Naimark-realization}
Every bounded positive operator-valued measure satisfies
\begin{equation}
 {
 \Sigma_\mu(z)
 =B_\mu(C_\mu-z)^{-1}B_\mu^*.}
 \label{eq:supp-Naimark-realization}
\end{equation}
Moreover
\begin{equation}
 {
 K_\mu
 =\overline{\Span}\{
 P_\mu(E)B_\mu^*\xi:
 E\in\mathfrak B((0,\infty)),\ \xi\in H\}.}
 \label{eq:supp-Naimark-minimal}
\end{equation}
Conversely, if $C\ge0$ is self-adjoint on $K$, $B:K\to H$ is bounded, and
\[
 K=\overline{\Span}\{E_C(E)B^*H\},
\]
then $\mu(E)=BE_C(E)B^*$ is a bounded positive operator-valued measure with
self-energy $B(C-z)^{-1}B^*$.  Two minimal realizations of the same
$\mu$ are related by a unique unitary $U$ satisfying
\begin{equation}
 {UC_1=C_2U,\qquad UB_1^*=B_2^*.}
 \label{eq:supp-Naimark-unitary}
\end{equation}
\end{theorem}

\begin{proof}
The spectral calculus gives
\[
 B_\mu(C_\mu-z)^{-1}B_\mu^*
 =V_\mu^*
 \left(\int\frac{1}{\lambda-z}\,\dd P_\mu(\lambda)\right)V_\mu
 =\int\frac{1}{\lambda-z}\,\dd\mu(\lambda).
\]
Every simple function is a finite sum of vectors
$P_\mu(E)V_\mu\xi$, proving \eqref{eq:supp-Naimark-minimal}.

For the converse, the spectral theorem gives positivity and countable
additivity of $BE_C(E)B^*$.  On the algebraic span of spectral-source vectors
define
\[
 U_0\left(\sum_jE_C(E_j)B^*\xi_j\right)
 =
 \sum_jP_\mu(E_j)V_\mu\xi_j.
\]
The squared norms on the two sides are both
\[
 \sum_{j,k}
 \langle\xi_j,BE_C(E_j\cap E_k)B^*\xi_k\rangle,
\]
so $U_0$ is well defined and isometric.  Source minimality makes its extension
unitary.  It intertwines every spectral projection, hence $C$, and sends
$B^*$ to $V_\mu$.  Applying the same construction to two realizations proves
existence and uniqueness of \eqref{eq:supp-Naimark-unitary}.
\end{proof}

\begin{theorem}[Static shadows, inverse moments, and positive Euclidean memory]
\label{thm:supp-dynamic-jets}
For $\operatorname{Im}z>0$,
\begin{equation}
 \operatorname{Im}\Sigma_\mu(z)
 =(\operatorname{Im}z)
  \int\frac{1}{|\lambda-z|^2}\,\dd\mu(\lambda)
 \succeq0.
 \label{eq:supp-Herglotz-sign}
\end{equation}
The finite static shadow
\[
 Q_0=\lim_{s\downarrow0}\Sigma_\mu(-s)
\]
exists exactly when the inverse first moment is bounded, and then
\begin{equation}
 {Q_0=\int\lambda^{-1}\,\dd\mu(\lambda).}
 \label{eq:supp-static-inverse-moment}
\end{equation}
If the corresponding inverse moment is finite, the $r$-th boundary jet is
\begin{equation}
 {
 \Sigma_{\mu,-}^{(r)}(0)
 =r!\int\lambda^{-(r+1)}\,\dd\mu(\lambda).}
 \label{eq:supp-dynamic-jet}
\end{equation}
The Euclidean memory kernel
\begin{equation}
 {
 K_\mu(t)=\int e^{-t\lambda}\,\dd\mu(\lambda)
 =B_\mu e^{-tC_\mu}B_\mu^*}
 \label{eq:supp-Euclidean-memory}
\end{equation}
is positive and
\[
 \Sigma_\mu(-s)
 =\int_0^\infty e^{-st}K_\mu(t)\,\dd t.
\]
\end{theorem}

\begin{proof}
The scalar identity
$\operatorname{Im}(\lambda-z)^{-1}
=(\operatorname{Im}z)|\lambda-z|^{-2}$ gives the sign after positive
operator integration.  Monotone convergence along $z=-s$ gives
\eqref{eq:supp-static-inverse-moment}.  Differentiating
$(\lambda-z)^{-1}$ gives the higher inverse moments.  The spectral calculus
gives the second equality in \eqref{eq:supp-Euclidean-memory}; Tonelli's
theorem applied to scalarizations gives the Laplace relation.
\end{proof}

\begin{remark}[Positive memory is not full OS positivity]
Positivity of $K_\mu(t)$ is a two-point positive memory statement.  It does not
by itself supply a Euclidean field law, a positive-time product algebra, or
Osterwalder--Schrader positivity for reflected words.
\end{remark}

\begin{theorem}[Determinacy on controlled support]
\label{thm:supp-memory-determinacy}
If $\supp\mu$ is compact, the complete nonnegative moment sequence
$M_n=\int\lambda^n\,\dd\mu(\lambda)$ determines $\mu$ uniquely.  The same
conclusion holds on unbounded support if
$\int e^{a\lambda}\,\dd\Tr\mu(\lambda)<\infty$ for some $a>0$.
\end{theorem}

\begin{proof}
On compact support, equality of all moments gives equality against all
polynomials; Stone--Weierstrass gives equality against every continuous
function, and polarization gives equality of the operator measures.

Under an exponential moment, each scalarized Laplace transform
$\int e^{-t\lambda}\,\dd\langle\xi,\mu\eta\rangle$ extends holomorphically to
a half-plane containing the origin.  Its Taylor coefficients are the moments.
Equality of all moments therefore gives equality of the Laplace transforms
near zero and hence throughout the connected half-plane.  Uniqueness of the
Laplace transform and polarization recover $\mu$.
\end{proof}

\begin{proposition}[All integer moments can miss the dynamic self-energy]
\label{cth:supp-moment-indeterminate}
For $|\varepsilon|\le1$ define
\begin{equation}
 {
 \dd\mu_\varepsilon(x)
 =
 \left[1+\varepsilon\sin(2\pi\log x)\right]
 \frac{e^{-(\log x)^2/2}}{x\sqrt{2\pi}}\,\dd x.}
 \label{eq:supp-lognormal-family}
\end{equation}
Then for every integer $n\in\mathbb Z$,
\begin{equation}
 {
 \int_0^\infty x^n\,\dd\mu_\varepsilon(x)=e^{n^2/2}}
 \label{eq:supp-lognormal-moments}
\end{equation}
independently of $\varepsilon$.  Nevertheless
$\varepsilon\ne\varepsilon'$ implies that the Stieltjes transforms differ for
some, hence an open set of, $z\in\mathbb C\setminus[0,\infty)$.  Consequently
all positive moments, all inverse integer moments, and every finite Hankel
panel can agree while the dynamic self-energies are inequivalent.
\end{proposition}

\begin{proof}
Set $y=\log x$.  The perturbation of the $n$-th moment is
\[
 \frac1{\sqrt{2\pi}}
 \int_\mathbb R
 e^{ny-y^2/2}\sin(2\pi y)\,\dd y
 =
 \operatorname{Im}
 \exp\!\left(\frac{(n+2\pi i)^2}{2}\right)
 =
 e^{(n^2-4\pi^2)/2}\sin(2\pi n)=0
\]
for every $n\in\mathbb Z$.  The unperturbed lognormal moment is
$e^{n^2/2}$.  Distinct $\varepsilon$ give distinct measures.

If a finite signed measure on $[0,\infty)$ had identically zero Stieltjes
transform on the negative real axis, the identity
\[
 \int\frac{1}{x+s}\,\dd\sigma(x)
 =
 \int_0^\infty e^{-st}
 \left(\int e^{-tx}\,\dd\sigma(x)\right)\dd t
\]
and two applications of uniqueness of the Laplace transform would force
$\sigma=0$.  Hence distinct measures have distinct Stieltjes transforms.
\end{proof}

\begin{theorem}[High-energy memory escape and instantaneous contact terms]
\label{thm:supp-contact-escape}
Let $\mu_n$ be positive tail measures supported in $[m,\infty)$ for one
$m>0$ and define
\[
 \dd\nu_n(\lambda)=\lambda^{-1}\dd\mu_n(\lambda),
 \qquad
 Q_n=\nu_n([m,\infty))=\Sigma_n(0).
\]
If $\sup_n\Tr Q_n<\infty$, every sequence has a subsequence and a positive
operator-valued measure on the one-point compactification of the form
\[
 \nu_\infty=\nu_{\rm fin}+Q_{\rm esc}\delta_\infty,
 \qquad Q_{\rm esc}\succeq0,
\]
such that, locally uniformly off $[m,\infty)$,
\begin{equation}
 {
 \Sigma_n(z)\longrightarrow
 Q_{\rm esc}
 +\int_{[m,\infty)}
  \frac{\lambda}{\lambda-z}\,\dd\nu_{\rm fin}(\lambda).}
 \label{eq:supp-contact-selfenergy}
\end{equation}
The Euclidean memory measures satisfy
\begin{equation}
 {
 K_n(t)\,\dd t
 \longrightarrow
 Q_{\rm esc}\delta_0+K_{\rm fin}(t)\,\dd t.}
 \label{eq:supp-contact-kernel}
\end{equation}
\end{theorem}

\begin{proof}
The matrix entries of $\nu_n$ are uniformly bounded finite measures on the
compact space $[m,\infty]$.  Weak-* compactness and polarization give a
positive operator-valued limit, which decomposes into finite and point-at-
infinity parts.  For $z\notin[m,\infty)$,
$g_z(\lambda)=\lambda/(\lambda-z)$ extends continuously to infinity with
$g_z(\infty)=1$, proving \eqref{eq:supp-contact-selfenergy}.
For $\varphi\in C_c([0,\infty))$,
\[
 h_\varphi(\lambda)
 =\int_0^\infty\varphi(t)\lambda e^{-\lambda t}\,\dd t
\]
extends continuously with $h_\varphi(\infty)=\varphi(0)$; integration against
$\nu_n$ gives \eqref{eq:supp-contact-kernel}.
\end{proof}

\begin{proposition}[The static triple does not identify escaped memory]
\label{cth:supp-contact-static}
For $Q\succeq0$ and $\lambda_n\to\infty$, let
\[
 \mu_n=\lambda_nQ\,\delta_{\lambda_n}.
\]
Then
\begin{equation}
 {
 \Sigma_n(0)=Q,\qquad
 \Sigma_n(z)=\frac{\lambda_n}{\lambda_n-z}Q\longrightarrow Q,}
 \label{eq:supp-contact-static}
\end{equation}
whereas $\lambda_ne^{-\lambda_nt}Q\to0$ for every $t>0$ and converges
distributionally to $Q\delta_0$.  The same static local correction is
therefore compatible with an intrinsic head term or with memory escaping to
infinite energy.
\end{proposition}

\begin{proof}
The resolvent identities are immediate, and the exponential kernels form the
standard positive approximate identity at the origin.
\end{proof}
\subsection{Global marked monodromy}

\begin{definition}[Marked spectral packet]
\label{def:supp-marked-spectral-packet}
A marked finite spectral packet is
\begin{equation}
 \mathsf S=(\A,\Hh,D,J,\gamma;\mathfrak m),
 \label{eq:supp-marked-packet}
\end{equation}
where $\mathfrak m$ denotes any declared auxiliary data that are required to be
preserved by isomorphisms, such as orientation, modular, calibration, frame, or
determinant-line markings.  Let
\begin{equation}
 G_{\mathsf S}=\operatorname{Aut}(\mathsf S)
 \label{eq:supp-marked-aut}
\end{equation}
be the compact group of unitary packet automorphisms preserving those
markings and the declared algebra automorphism convention.
\end{definition}

\begin{theorem}[Marked spectral-bundle descent]
\label{thm:supp-marked-bundle}
Let $\Gamma$ be connected, let each vertex packet be marked-isomorphic to
$\mathsf S$, and let each oriented edge carry a marked unitary isomorphism
$U_e$ with $U_{\bar e}=U_e^{-1}$.  Then:
\begin{enumerate}[label=\textup{(B\arabic*)},leftmargin=3em]
\item after choosing a base vertex and spanning tree, chord holonomies define a
      homomorphism
      \begin{equation}
       \rho_U:\pi_1(\Gamma)\longrightarrow G_{\mathsf S};
       \label{eq:supp-monodromy}
      \end{equation}
\item changing vertex frames conjugates $\rho_U$ by one element of
      $G_{\mathsf S}$;
\item one global marked trivialization exists if and only if $\rho_U$ is
      trivial;
\item in general the isomorphism class of the flat marked spectral bundle is
      \begin{equation}
       {
       [\rho_U]\in
       \operatorname{Hom}(\pi_1(\Gamma),G_{\mathsf S})/G_{\mathsf S}.}
       \label{eq:supp-monodromy-class}
      \end{equation}
\end{enumerate}
Since $\pi_1(\Gamma)$ is free of rank $b_1(\Gamma)$, the bundle is determined
by $b_1(\Gamma)$ elements of $G_{\mathsf S}$ modulo simultaneous conjugation.
\end{theorem}

\begin{proof}
Use the edge maps along a spanning tree to identify every vertex packet with
the base packet.  Each chord gives the ordered product around its fundamental
cycle, an element of $G_{\mathsf S}$.  These cycle values extend uniquely to a
homomorphism from the free fundamental group.  Changing the base frame
conjugates every cycle value by the same automorphism.  Trivial cycle values
make the tree identification path independent and hence give one global marked
frame.  Conversely, a global frame makes every cycle product telescope.  This
is precisely the clutching classification
\eqref{eq:supp-monodromy-class} on a graph.
\end{proof}

\begin{corollary}[Nontrivial monodromy is a twisted geometry]
\label{cor:supp-twisted-bundle}
If the local packets are marked-isomorphic and the edge identifications satisfy the declared compatibility relations while $[\rho_U]$ is nontrivial, the resulting object is a flat twisted spectral geometry rather than a globally trivial marked packet.
\end{corollary}

\begin{theorem}[Line characters are incomplete shadows of frame holonomy]
\label{thm:supp-character-shadow}
Let
$\chi_1,\ldots,\chi_m:G_{\mathsf S}\to U(1)$ be determinant, charge, or other
complex line characters and let
$\epsilon_1,\ldots,\epsilon_r:G_{\mathsf S}\to\{\pm1\}$ be real Pfaffian-sign
characters.  Trivial frame holonomy implies triviality of all induced line
holonomies.  The converse holds exactly when
\begin{equation}
 {
 \bigcap_j\Ker\chi_j\cap\bigcap_k\Ker\epsilon_k=\{1\}.}
 \label{eq:supp-character-faithfulness}
\end{equation}
\end{theorem}

\begin{proof}
Every character sends the identity to the identity.  If
\eqref{eq:supp-character-faithfulness} holds and all character holonomies are
trivial, every cycle value of $\rho_U$ lies in the displayed intersection and
is therefore the identity.  Conversely, if the intersection contains
$g\ne1$, assign $g$ as monodromy on a one-cycle graph.  Every listed character
is trivial although the marked bundle is not.
\end{proof}

\begin{proposition}[A trivial determinant line need not trivialize the packet]
\label{cth:supp-det-no-frame}
Let a marked packet have scalar algebra, zero Dirac operator, and a
multiplicity-two Hilbert carrier so that its marked automorphism group contains
$SU(2)$.  On a one-cycle graph choose monodromy
\begin{equation}
 g=\begin{pmatrix}e^{i\theta}&0\\0&e^{-i\theta}\end{pmatrix},
 \qquad \theta\notin2\pi\Z.
 \label{eq:supp-SU2-monodromy}
\end{equation}
Then the determinant-line holonomy is one although the marked Hilbert bundle
has nontrivial monodromy.
\end{proposition}

\begin{proof}
Every $SU(2)$ element has determinant one, while the displayed matrix is not
the identity.  Apply \Cref{thm:supp-marked-bundle,thm:supp-character-shadow}.
\end{proof}
\subsection{Coarse graining of the principal fibres}

Let $\partial_G:\mathbb R^E\to\mathbb R^V$ be the oriented incidence map and
$\mathcal Z(G)=\Ker\partial_G$.  A vertex quotient $q:V\twoheadrightarrow\bar V$
induces maps $Q_V$ and $Q_E$ by summing vertex and oriented edge coefficients over
fibres.

\begin{theorem}[Current-chain functor]
\label{thm:supp-current-chain}
The incidence maps satisfy
\begin{equation}
{\partial_{\bar G}Q_E=Q_V\partial_G,}
\label{eq:supp-current-chain}
\end{equation}
so $Q_E\mathcal Z(G)\subseteq\mathcal Z(\bar G)$.  If
$\bar c_{\bar e}=\sum_{e\mapsto\bar e}c_e$, then
$Q_E\Curr(G,c)\subseteq\Curr(\bar G,\bar c)$.  When the vertex fibres are
connected and the coarse graph retains every inter-fibre edge class, the induced
map on cycle spaces is onto and
\begin{equation}
0\longrightarrow\Ker(Q_E|_{\mathcal Z(G)})
\longrightarrow\mathcal Z(G)
\xrightarrow{Q_E}\mathcal Z(\bar G)
\longrightarrow0,
\label{eq:supp-cycle-exact}
\end{equation}
with
\begin{equation}
\dim\Ker(Q_E|_{\mathcal Z(G)})=b_1(G)-b_1(\bar G).
\label{eq:supp-cycle-kernel-dim}
\end{equation}
The stationary entropy production satisfies
\begin{equation}
\sigma_{\bar G}(Q_Ej)\le\sigma_G(j).
\label{eq:supp-entropy-contraction}
\end{equation}
\end{theorem}

\begin{proof}
For a coarse vertex, the divergence of the summed inter-fibre current equals the
sum of fine divergences; internal fine edges cancel in pairs.  This is exactly
\eqref{eq:supp-current-chain}.  The conductance inequality
$|(Q_Ej)_{\bar e}|\le\sum_{e\mapsto\bar e}|j_e|<\bar c_{\bar e}$ gives the fibre
map.  Connected fibres allow every coarse cycle to be lifted and completed inside
fibres, proving surjectivity.  The dimension formula follows from the Betti-number
formula for connected graphs.  Finally, on every coarse edge the log-sum inequality
for the aggregated forward and reverse fluxes gives entropy contraction; summing
over coarse edges yields \eqref{eq:supp-entropy-contraction}.
\end{proof}

For an effective operator $A$ and positive memory measure $\mu$, write
\begin{equation}
F_{A,\mu}(z)=A-z-\Sigma_\mu(z),
\qquad
Q_{>\Lambda}=\int_{(\Lambda,\infty)}\lambda^{-1}\,\dd\mu(\lambda),
\label{eq:supp-memory-head}
\end{equation}
and define
\begin{equation}
\RG_\Lambda(A,\mu)
=\left(A-Q_{>\Lambda},\mu|_{(0,\Lambda]}\right).
\label{eq:supp-memory-RG}
\end{equation}

\begin{theorem}[Exact memory RG semigroup]
\label{thm:supp-memory-RG}
If $(A_\Lambda,\mu_\Lambda)=\RG_\Lambda(A,\mu)$, then
\begin{equation}
{F_{A_\Lambda,\mu_\Lambda}(0)=F_{A,\mu}(0).}
\label{eq:supp-memory-static-preserved}
\end{equation}
For $|z|\le r<\Lambda$,
\begin{equation}
{
\|F_{A,\mu}(z)-F_{A_\Lambda,\mu_\Lambda}(z)\|
\le\frac r{\Lambda-r}\|Q_{>\Lambda}\|.}
\label{eq:supp-memory-screen-error}
\end{equation}
Moreover, if $0<\Lambda_2\le\Lambda_1$,
\begin{equation}
\RG_{\Lambda_2}\circ\RG_{\Lambda_1}=\RG_{\Lambda_2}.
\label{eq:supp-memory-semigroup}
\end{equation}
\end{theorem}

\begin{proof}
Split the Stieltjes integral at $\Lambda$.  At $z=0$ the removed high-energy part is
exactly $Q_{>\Lambda}$ and is canceled by the head correction.  Away from zero,
\[
\frac1{\lambda-z}-\frac1\lambda
=\frac z{\lambda(\lambda-z)},
\]
whose norm is at most
$r[(\Lambda-r)\lambda]^{-1}$ for $\lambda>\Lambda$.  Integrating gives
\eqref{eq:supp-memory-screen-error}.  Two successive cutoffs simply split the same
measure into disjoint shells; their inverse-moment head corrections add, proving the
semigroup law.
\end{proof}

Let $\rho:\pi_1(\Gamma,v_0)\to G_{\mathsf S}$ be marked monodromy and let a graph
contraction $q:\Gamma\to\bar\Gamma$ induce a surjection $q_*$ on fundamental
groups.

\begin{theorem}[Exact marked-monodromy descent]
\label{thm:supp-monodromy-descent}
There is a unique $\bar\rho:\pi_1(\bar\Gamma)\to G_{\mathsf S}$ satisfying
\begin{equation}
\rho=\bar\rho\circ q_*
\label{eq:supp-monodromy-factor}
\end{equation}
if and only if
\begin{equation}
{\Ker q_*\subseteq\Ker\rho.}
\label{eq:supp-monodromy-kernel}
\end{equation}
The conjugacy class $[\bar\rho]$ is gauge independent and functorial under further
contractions.  A determinant or Pfaffian character may satisfy its own descended
kernel condition even when \eqref{eq:supp-monodromy-kernel} fails, so line descent
does not imply descent of the full marked packet.
\end{theorem}

\begin{proof}
If \eqref{eq:supp-monodromy-factor} holds, every element of $\Ker q_*$ lies in
$\Ker\rho$.  Conversely, if \eqref{eq:supp-monodromy-kernel} holds, define
$\bar\rho(q_*\gamma)=\rho(\gamma)$.  The kernel condition makes the definition
independent of the chosen lift.  Surjectivity of $q_*$ gives uniqueness and the
quotient-group universal property gives functoriality.  The character statement
follows by applying the same argument to a nonfaithful one-dimensional
representation of $G_{\mathsf S}$.
\end{proof}

\begin{theorem}[Functorial principal universality packet]
\label{thm:supp-fibre-RG}
Collect
\begin{equation}
\mathfrak U(G,c;A,\mu;\rho)
=\bigl(\Curr(G,c),(A,\mu),[\rho]\bigr).
\label{eq:supp-U-packet}
\end{equation}
For a connected-fibre graph quotient, a memory threshold $\Lambda$, and a compatible
isometric effective-space compression $V$, define
\begin{equation}
{
\mathfrak U\longmapsto
\left(
Q_E\Curr(G,c),
\bigl(V^*(A-Q_{>\Lambda})V,
      V^*\mu|_{(0,\Lambda]}V\bigr),
[\bar\rho]
\right).}
\label{eq:supp-U-map}
\end{equation}
When the monodromy kernel condition holds, this map is canonical and composes
strictly for composable graph quotients, nested memory thresholds and compatible
head maps.  Along it, entropy production cannot increase; the zero-energy Schur complement is preserved before effective-space compression; minimal dilation rank and
Pontryagin negative-square index cannot increase; and marked monodromy descends
exactly when collapsed cycles lie in its kernel.
\end{theorem}

\begin{proof}
The current and entropy statements are \Cref{thm:supp-current-chain}; the memory
statement and exact composition are \Cref{thm:supp-memory-RG}; the monodromy part is
\Cref{thm:supp-monodromy-descent}.  Compression of a positive minimal memory dilation
carrier cannot increase source rank.  For a generalized-Nevanlinna kernel, head
compression replaces each finite kernel Gram by a congruence, which cannot increase
the number of negative eigenvalues.  Each component is uniquely determined by the
corresponding quotient data, so the product map is canonical and composes
componentwise.
\end{proof}

\section{Analytic details for compact-spin spectralization}
\label{app:compact-spin}
% =============================================================================

The periodic $A_3$ construction in Section~\ref{sec:metric} supplies the explicit metric model.  This appendix records the analytic estimates used in the compact-spin convergence criterion.  The chartwise construction is verified below; the global theorem additionally assumes the atlas stability conditions of Definition~\ref{def:stable-spin-atlas}.

\subsection{A positive Renewal stencil for an arbitrary local inverse metric}

Let $d\ge2$ and put
\begin{equation}
\mathcal R_d^+
=\{e_i:1\le i\le d\}
\cup\{e_i+e_j,e_i-e_j:1\le i<j\le d\}.
\label{eq:supp-general-stencil}
\end{equation}
There are exactly $d^2$ unoriented directions.

\begin{lemma}[Positive $d^2$-direction decomposition]
\label{lem:supp-positive-decomposition}
Let $A=(A_{ij})\in\operatorname{Sym}_d(\mathbb R)$ and suppose that for some
$\eta>0$,
\begin{equation}
|A_{ij}|<\eta\quad(i\ne j),
\qquad
A_{ii}>(d-1)\eta.
\label{eq:supp-stencil-window}
\end{equation}
Define
\begin{equation}
w_i=A_{ii}-(d-1)\eta,
\qquad
w_{ij}^{\pm}=\frac12(\eta\pm A_{ij}).
\label{eq:supp-stencil-weights}
\end{equation}
Then every coefficient is strictly positive and
\begin{equation}
{
A=\sum_iw_ie_ie_i^{\mathsf T}
+\sum_{i<j}\left[
 w_{ij}^+(e_i+e_j)(e_i+e_j)^{\mathsf T}
+w_{ij}^-(e_i-e_j)(e_i-e_j)^{\mathsf T}
\right].}
\label{eq:supp-positive-decomposition}
\end{equation}
\end{lemma}

\begin{proof}
Positivity follows immediately from \eqref{eq:supp-stencil-window}.  In the
right-hand side of \eqref{eq:supp-positive-decomposition}, the $(i,j)$ off-diagonal
entry is $w_{ij}^+-w_{ij}^-=A_{ij}$.  Every pair containing $i$ contributes
$w_{ij}^++w_{ij}^-=\eta$ to the $i$th diagonal, so the diagonal coefficient is
$w_i+(d-1)\eta=A_{ii}$.
\end{proof}

A linear coordinate normalization makes $g^{-1}$ equal to $I_d$ at the centre of a
chart.  By continuity, after shrinking the chart, the inverse metric remains inside
one window \eqref{eq:supp-stencil-window}.  Compactness gives a finite good cover by
such charts.

Let $\rho=\sqrt{\det g}$ and write the decomposition coefficients of $g^{-1}$ as
$w_r(x)$, $r\in\mathcal R_d^+$.  Put $a_r(x)=\rho(x)w_r(x)$.  On
$Q_h=h\mathbb Z^d/\mathbb Z^d$ take
\begin{equation}
m_h(x)=h^d\rho(x),
\qquad
c_h(x,r)=\frac{h^{d-2}}4\bigl(a_r(x)+a_r(x+hr)\bigr).
\label{eq:supp-general-conductance}
\end{equation}
The two directed rates are the conductance divided by the corresponding endpoint
mass.

\begin{theorem}[Reversible coframe Renewal process]
\label{thm:supp-general-renewal-process}
The packet \eqref{eq:supp-general-conductance} is reversible with stationary mass
$m_h$.  Its predictable coordinate bracket is
\begin{equation}
{
\mathcal B_h(x)
=\frac1{4\rho(x)}
\sum_{r\in\mathcal R_d^+}
\bigl[a_r(x+hr)+2a_r(x)+a_r(x-hr)\bigr]rr^{\mathsf T},}
\label{eq:supp-general-bracket}
\end{equation}
and satisfies
\begin{equation}
\|\mathcal B_h-g^{-1}\|_{L^\infty}\le Ch^2.
\label{eq:supp-general-bracket-error}
\end{equation}
For every $f\in C^4$,
\begin{equation}
{\mathcal L_hf=\frac12\Delta_gf+O(h^2\|f\|_{C^4})}
\label{eq:supp-general-generator}
\end{equation}
uniformly on the chart, and the Dirichlet forms converge to
\begin{equation}
-\langle f_h,\mathcal L_hf_h\rangle_{m_h}
\longrightarrow
\frac12\int g^{ij}\partial_if\,\partial_j\overline f\,\rho\,\dd x.
\label{eq:supp-general-form-limit}
\end{equation}
\end{theorem}

\begin{proof}
Symmetry of $c_h$ gives detailed balance.  The bracket follows by summing the two
orientations of every direction through a vertex.  Taylor expansion of
$a_r(x\pm hr)$ cancels odd terms and, together with
\eqref{eq:supp-positive-decomposition}, gives
\eqref{eq:supp-general-bracket-error}.  Expanding
$f(x\pm hr)$ to fourth order, the first-order terms cancel and the second-order
terms contract with the same tensor decomposition.  The derivative of the physical
density and of the conductances produces the divergence-form first-order part of
$\Delta_g$, yielding \eqref{eq:supp-general-generator}.  Summation by parts gives
the finite Dirichlet form; the coefficient convergence and Riemann sums give
\eqref{eq:supp-general-form-limit}.
\end{proof}

\begin{proposition}[The metric bracket does not select the spin marking]
\label{cth:supp-metric-no-coframe}
The positive bracket determines $g^{-1}$ but not an oriented orthonormal coframe, a
spin lift, or a global spin structure.  A smooth rotation
$R(x)\in\mathrm{SO}(d)$ changes the coframe and spin connection without changing
$g^{-1}$ or the scalar Renewal generator, and a nontrivial $\mathbb Z_2$ spin
holonomy is invisible to every local metric bracket.
\end{proposition}

\subsection{Density-symmetric spin Dirac and finite Wilson packet}

Let $S_d$ be a complex spin module, let $\gamma_a$ be Hermitian Clifford matrices,
and let $E_a{}^j$ be the inverse oriented orthonormal coframe.  Set
\begin{equation}
c^j=E_a{}^j\gamma_a,
\qquad
c^jc^k+c^kc^j=2g^{jk}I,
\label{eq:supp-clifford-coeff}
\end{equation}
and let $\Omega_j$ be the Levi--Civita spin connection.  On
$L^2(Q,S_d;\rho\dd x)$,
\begin{equation}
D_g=-i\sum_jc^j(\partial_j+\Omega_j).
\label{eq:supp-geometric-dirac}
\end{equation}

\begin{lemma}[Exact density-symmetric Dirac identity]
\label{lem:supp-density-symmetric}
With $U_\rho\psi=\rho^{1/2}\psi$ and
$\nabla_j=\partial_j+\Omega_j$,
\begin{equation}
{
U_\rho D_gU_\rho^{-1}
=-\frac{i}{2}\sum_{j=1}^d
\bigl(M_{c^j}\nabla_j+\nabla_jM_{c^j}\bigr).}
\label{eq:supp-density-symmetric}
\end{equation}
\end{lemma}

\begin{proof}
Metric compatibility of Clifford multiplication gives
\[
\partial_jc^k+[\Omega_j,c^k]+\Gamma^k{}_{j\ell}c^\ell=0.
\]
Setting $k=j$ and using
$\Gamma^j{}_{j\ell}=\partial_\ell\log\rho$ gives
\[
\partial_jc^j+[\Omega_j,c^j]+c^j\partial_j\log\rho=0.
\]
Conjugating \eqref{eq:supp-geometric-dirac} by $\rho^{1/2}$ and applying this
contracted identity yields \eqref{eq:supp-density-symmetric}.
\end{proof}

Introduce one fixed two-dimensional normal factor and put
\begin{equation}
\widehat S_d=S_d\otimes\mathbb C^2,
\quad
\widehat c^j=c^j\otimes\sigma_1,
\quad
\Gamma_\perp=I\otimes\sigma_2,
\quad
\widehat\gamma=I\otimes\sigma_3.
\label{eq:supp-doubled-clifford}
\end{equation}
Let $T_{j,h}^{\Omega}$ be exact inverse spin parallel transport followed by a lattice
shift and define
\begin{equation}
P_{j,h}^{\Omega}
=\frac{T_{j,h}^{\Omega}-(T_{j,h}^{\Omega})^*}{2ih},
\qquad
W_h^{\Omega}
=\frac1{2h}\sum_{j=1}^d
\bigl(2I-T_{j,h}^{\Omega}-(T_{j,h}^{\Omega})^*\bigr).
\label{eq:supp-covariant-differences}
\end{equation}
The local doubled Wilson operator is the density-conjugated version of
\begin{equation}
\widetilde D_{g,h}^{\rm W}
=\frac12\sum_j
\bigl(M_{\widehat c_h^j}P_{j,h}^{\Omega}
      +P_{j,h}^{\Omega}M_{\widehat c_h^j}\bigr)
+\varpi\Gamma_\perp W_h^{\Omega}.
\label{eq:supp-general-Wilson}
\end{equation}

For the local estimates, extend the coefficients smoothly to a fixed periodic
coordinate box, preserving uniform ellipticity, and use the density-conjugated
$\ell^2$ norm with cell weight $h^d$.  Such an extension is an auxiliary local
calculation, not a boundary condition on a chart of $M$.  Let $S_{j,h}$ be the
ordinary coordinate shift.  In a smooth spin frame,
$T_{j,h}^{\Omega}-S_{j,h}=O(h)$ in operator norm.  Define
\begin{equation}
 \|u\|_{1,h}^2=\|u\|_h^2+
       \sum_j\left\|h^{-1}(S_{j,h}-I)u\right\|_h^2.
 \label{eq:discrete-H1-norm}
\end{equation}
The analogous covariant norm is uniformly equivalent to this one, since the
connection coefficients are bounded.  All constants below may depend on the
fixed atlas and finitely many coefficient bounds, but not on $h$.

\begin{theorem}[Frozen symbol, unique species, and uniform graph estimate]
\label{thm:supp-general-Wilson-ellipticity}
For fixed $\varpi>0$, the frozen principal symbol of
\eqref{eq:supp-general-Wilson} is
\begin{equation}
 q_{h,x_0}(\theta)=\frac1h\left[
 \sum_j\widehat c^j(x_0)\sin\theta_j
 +\varpi\Gamma_\perp\sum_j(1-\cos\theta_j)\right].
 \label{eq:supp-general-frozen-symbol}
\end{equation}
It satisfies
\begin{equation}
 q_{h,x_0}(\theta)^2=\frac1{h^2}\left[
 g^{jk}(x_0)\sin\theta_j\sin\theta_k
 +\varpi^2\Bigl(\sum_j(1-\cos\theta_j)\Bigr)^2\right]I.
 \label{eq:supp-general-frozen-square}
\end{equation}
Its only zero is $\theta=0$ on the Brillouin torus.  On the fixed smooth
periodic extension, and for sections localized in the interior of the original
chart, there is a cutoff-independent bound
\begin{equation}
 \|u\|_{1,h}\le C\bigl(\|\widetilde D_{g,h}^{\rm W}u\|_h+\|u\|_h\bigr).
 \label{eq:supp-general-Garding}
\end{equation}
The symmetric placement of coefficients makes the finite operator self-adjoint;
the fixed doubled Clifford convention supplies the declared real and grading
structures.
\end{theorem}

\begin{proof}
The Clifford relations eliminate the mixed terms in the square.  Put
$s(\theta)=\sum_j(1-\cos\theta_j)$ and
$\ell_h(\theta)=2h^{-2}s(\theta)$.  Uniform ellipticity and the elementary
identity
\[
 \sum_j\sin^2\theta_j+s(\theta)^2
 =2s(\theta)+2\sum_{i<j}(1-\cos\theta_i)(1-\cos\theta_j)
 \ge2s(\theta)
\]
give constants $0<c\le C<\infty$ such that
\begin{equation}
 c\ell_h(\theta)I\preceq q_{h,x_0}(\theta)^2
            \preceq C\ell_h(\theta)I.
 \label{eq:frozen-uniform-ellipticity}
\end{equation}
For the upper bound use $s\le2d$ and $\sum_j\sin^2\theta_j\le2s$.
The lower bound controls all forward differences, including modes near the
nonzero Brillouin corners.  Fourier transformation proves the frozen graph
estimate.

To pass to variable coefficients, choose a partition with supports of fixed
small diameter $r$, independent of $h$.  On one support the difference of the
variable principal part and its frozen value is bounded by
$C\omega_c(r)\|u\|_{1,h}+C_r\|u\|_h$, where $\omega_c$ is the modulus of
continuity of the coframe coefficients.  Derivatives of coefficients and the
replacement of ordinary shifts by smooth spin links contribute only the
second term.  Choose $r$ once so that the first term is absorbed by the frozen
estimate.  For every fixed smooth cutoff $\chi$,
\begin{equation}
 \|[\widetilde D_{g,h}^{\rm W},M_{\chi,h}]\|
          \le C\|\nabla\chi\|_\infty.
 \label{eq:local-cutoff-commutator}
\end{equation}
Indeed, a shift commutator contains $\chi(x+he_j)-\chi(x)=O(h)$,
which cancels the $h^{-1}$ factor in both the Dirac and Wilson differences.
Apply the local estimate to $\chi u$, sum the squared bounds over a fixed
finite partition with $\sum\chi^2=1$, and use the same discrete product rule
for the forward differences in \eqref{eq:discrete-H1-norm}.  This proves
\eqref{eq:supp-general-Garding}.  The fixed localization scale is important:
no cutoff derivative diverging as $h\to0$ is absorbed into a uniform constant.
\end{proof}

Let $\mathcal J_h^0$ be isometric piecewise-constant reconstruction and take
$S_h=(\mathcal J_h^0)^*$, the cell-average sampling map in the chosen density
and spin frame.  This avoids a dimension-dependent point-evaluation assumption
on Sobolev functions.  A piecewise-affine reconstruction $I_h$ on a fixed
shape-regular refinement of the coordinate cells satisfies
\begin{equation}
 \|I_hu\|_{H^1}\le C\|u\|_{1,h},\qquad
 \|I_hu-\mathcal J_h^0u\|_{L^2}\le Ch\|u\|_{1,h}.
 \label{eq:local-reconstruction-estimates}
\end{equation}
These inequalities follow cell by cell from the affine interpolation formula
and scaling to a unit cell.  Thus it is $I_hu$, not the discontinuous
piecewise-constant reconstruction itself, that has a uniform $H^1$ bound.

\begin{lemma}[Smooth-core consistency and the Wilson term]
\label{lem:supp-general-core}
For every smooth doubled spinor supported in the coordinate interior,
\begin{equation}
 \left\|\mathcal J_h^0\widetilde D_{g,h}^{\rm W}S_h\psi
       -U_\rho\widehat D_gU_\rho^{-1}\psi\right\|_{L^2}
       \le Ch\|\psi\|_{H^3}.
 \label{eq:supp-general-core}
\end{equation}
In particular, on this core,
\begin{equation}
 \|W_h^\Omega S_h\psi\|_h\le Ch\|\psi\|_{H^2}.
 \label{eq:Wilson-core-smallness}
\end{equation}
This is a core estimate, not an operator-norm smallness assertion on the full
lattice space.
\end{lemma}

\begin{proof}
Taylor's formula with integral remainder, applied to spin-parallel translated
cell averages, gives the covariant centered derivative $-i\nabla_j$ and the
second-difference expansion $W_h^\Omega=-(h/2)\sum_j\nabla_j^2$ on a smooth
core, with the corresponding Sobolev remainders.  Cell averaging and
reconstruction add $O(h)$ errors controlled by one further derivative.
Symmetric multiplication by the smooth coefficients and
Lemma~\ref{lem:supp-density-symmetric} give
\eqref{eq:supp-general-core}; applying the second-difference estimate alone
gives \eqref{eq:Wilson-core-smallness}.
\end{proof}

The two roles of the Wilson term must be distinguished.  On smooth sections it
vanishes as $h\to0$, whereas at nonzero lattice corners its size is of order
$h^{-1}$ and removes the spurious low-energy species.  It is therefore not a
bounded perturbation tending to zero in operator norm.  This distinction is
also central in rigorous continuum-limit studies of lattice Dirac operators
\cite{Nakamura2024}.  Moreover, for any bounded sequence $u_h$ and fixed smooth
$\psi$, self-adjointness gives
\[
 |\langle W_h^\Omega u_h,S_h\psi\rangle_h|
 =|\langle u_h,W_h^\Omega S_h\psi\rangle_h|
 \le Ch\|u_h\|_h\|\psi\|_{H^2}.
\]
Hence the regularizer leaves no additional distributional term in the limiting
Dirac equation.  On a closed manifold the intended self-adjoint realization is
the closure of $D_g$ on smooth spinors, with domain $H^1(M;S)$
\cite{lawson-michelsohn}; no chart-boundary domain is introduced.

\subsection{Atlas compatibility and the global resolvent argument}
\label{subsec:atlas-compatibility}

Local consistency and a positive overlap Gram do not, by themselves, imply
global norm-resolvent convergence.  The following formulation states the
required global approximation properties explicitly, so that the analytic
conclusion is separate from a choice of mesh and overlap realization.

\begin{definition}[Compatible stable spin discretization]
\label{def:stable-spin-atlas}
Let $\widehat D=D_{M,\mathfrak s}\otimes\sigma_1$ on
$\mathscr H=L^2(M;S\otimes\C^2)$, where $M$ is closed.  A compatible stable
spin discretization is a sequence of finite Hilbert spaces $\Hh_h$, self-adjoint
operators $D_h$, and isometries $W_h:\Hh_h\to\mathscr H$ satisfying:
\begin{enumerate}[label=\textup{(A\arabic*)},leftmargin=3em]
\item $W_hW_h^*\to I$ strongly.  There are discrete first-order norms
$\|\cdot\|_{1,h}$ and reconstructions $I_h:\Hh_h\to H^1(M;S\otimes\C^2)$
with
\begin{equation}
 \|I_hu\|_{H^1}\le C\|u\|_{1,h},\qquad
 \|I_hu-W_hu\|\le\varepsilon_h\|u\|_{1,h},\quad
 \varepsilon_h\longrightarrow0.
 \label{eq:atlas-reconstruction}
\end{equation}
\item The assembled, rather than merely chartwise, operators satisfy
\begin{equation}
 \|u\|_{1,h}\le C(\|D_hu\|_h+\|u\|_h).
 \label{eq:atlas-global-stability}
\end{equation}
\item For every smooth doubled spinor $\psi$ there are samples $S_h\psi$ with
\begin{equation}
 W_hS_h\psi\longrightarrow\psi,\qquad
 W_hD_hS_h\psi\longrightarrow\widehat D\psi.
 \label{eq:atlas-global-core}
\end{equation}
The global samples use the transition functions of the fixed spin bundle on
shared degrees of freedom, including coordinate and density changes.  When
real and grading structures are asserted, their finite actions intertwine
these reconstructions and converge to the declared continuum actions.
\end{enumerate}
\end{definition}

Here is the cancellation that a proposed atlas assembly must respect.  Choose
fixed smooth real cutoffs $\chi_\alpha$ supported in chart interiors with
$\sum_\alpha\chi_\alpha^2=1$.  After transporting local spinors and their
density factors to the same bundle, the continuum first-order identity is
\begin{equation}
 \sum_\alpha\chi_\alpha\widehat D\chi_\alpha=\widehat D,
 \qquad
 \sum_\alpha\chi_\alpha[\widehat D,\chi_\alpha]
       =\tfrac12[\widehat D,\sum_\alpha\chi_\alpha^2]=0.
 \label{eq:atlas-cutoff-cancellation}
\end{equation}
Thus the apparent order-one derivatives of the cutoffs cancel; one must not
assert that they individually tend to zero.  For a finite-range Wilson
operator on a common coordinate mesh the analogous localization identity is
\begin{equation}
 \sum_\alpha M_{\chi_\alpha,h}W_h^\Omega M_{\chi_\alpha,h}-W_h^\Omega
 =\tfrac12\sum_\alpha
 [M_{\chi_\alpha,h},[W_h^\Omega,M_{\chi_\alpha,h}]],
 \label{eq:Wilson-double-commutator}
\end{equation}
and its norm is at most $Ch\sum_\alpha\|\nabla\chi_\alpha\|_\infty^2$.
Each matrix entry contains two cutoff differences and only one factor $h^{-1}$.
Fixed collars keep all such finite-range stencils away from auxiliary chart
boundaries for sufficiently small $h$.

On a shared geometric link $x\to y$, spin transition matrices satisfy the
exact covariance identity
\[
 U_{\alpha\beta}(x)\,T^\beta_{xy}\,U_{\alpha\beta}(y)^{-1}
       =T^\alpha_{xy}.
\]
It removes frame mismatch on that link.  It does not identify different
coordinate-grid links or supply a stable interpolation between unrelated grids.
For independently meshed charts, such interpolation and its first-order error
bounds must be verified in \eqref{eq:atlas-reconstruction}--
\eqref{eq:atlas-global-core}.  In particular, positivity of a Gram matrix is
not a substitute for the assembled graph estimate
\eqref{eq:atlas-global-stability}.  The local Wilson estimates above establish
the local ingredients; the existence of a global assembly satisfying all these
properties is an additional construction requirement.

\begin{theorem}[Norm-resolvent convergence for compatible stable atlases]
\label{thm:supp-compact-spin-resolvent}
For a compatible stable spin discretization in the sense of
Definition~\ref{def:stable-spin-atlas}, and every $z\notin\R$,
\begin{equation}
 \left\|W_h(D_h-z)^{-1}W_h^*-(\widehat D-z)^{-1}\right\|
       \longrightarrow0.
 \label{eq:supp-global-norm-resolvent}
\end{equation}
The same argument on a fixed smooth periodic coordinate extension gives
\begin{equation}
 \left\|\mathcal J_h^0(D_{g,h}^{\rm W}-z)^{-1}(\mathcal J_h^0)^*
       -(\widehat D_g-z)^{-1}\right\|\longrightarrow0.
 \label{eq:supp-local-norm-resolvent}
\end{equation}
This local statement is about that self-adjoint periodic extension, not an
unspecified self-adjoint operator on an open chart.  Isolated bounded-energy
spectral projections converge, and compatible real and grading structures
are retained as specified in Definition~\ref{def:stable-spin-atlas}.
\end{theorem}

\begin{proof}
Write $R_h(z)=W_h(D_h-z)^{-1}W_h^*$ and
$R(z)=(\widehat D-z)^{-1}$.  Self-adjointness gives
\begin{equation}
 \|(D_h-z)^{-1}\|\le|\operatorname{Im}z|^{-1},\qquad
 \|D_h(D_h-z)^{-1}\|\le1+|z|\,|\operatorname{Im}z|^{-1}.
 \label{eq:uniform-resolvent-bounds}
\end{equation}
For $\|f\|\le1$ put $u_h=(D_h-z)^{-1}W_h^*f$.  The global graph estimate
bounds $\|u_h\|_{1,h}$ uniformly.  Therefore $I_hu_h$ lies in a bounded
$H^1$ set, while $\|W_hu_h-I_hu_h\|\to0$ uniformly on this resolvent unit
ball.  Rellich compactness proves collective compactness of the embedded
resolvents along the refining sequence.  This step uses the reconstruction
estimate and does not assign an $H^1$ norm to a piecewise-constant function.

For fixed $f$, let $W_hu_h\to u$ along a subsequence.  Test the discrete
resolvent equation against $S_h\psi$, with $\psi$ any smooth doubled spinor.
Self-adjointness and \eqref{eq:atlas-global-core} pass the equation to
$(\widehat D-z)u=f$ in the distributional sense.  Since smooth spinors are a
core for the self-adjoint operator $\widehat D$, this places $u$ in its
adjoint domain, which is $H^1$, and gives $u=R(z)f$.  Every convergent
subsequence has this limit, so $R_h(z)\to R(z)$ strongly.  Applying the same
argument at $\bar z$ also gives strong convergence of the adjoints.

For clarity, this last fact is needed for the upgrade to norm convergence.
If that convergence failed, choose unit vectors $f_h$ with
$\|(R_h(z)-R(z))f_h\|\ge\epsilon$.  Passing to subsequences, let
$f_h\rightharpoonup f$ and use collective compactness to make $R_h(z)f_h$
converge in norm.  For each fixed $v$, strong adjoint convergence implies
\[
 \langle v,R_h(z)f_h\rangle
 =\langle R_h(\bar z)v,f_h\rangle
 \longrightarrow\langle R(\bar z)v,f\rangle.
\]
Hence its norm limit is $R(z)f$.  Compactness of $R(z)$ also gives
$R(z)f_h\to R(z)f$, a contradiction.  This proves
\eqref{eq:supp-global-norm-resolvent}.  On the periodic extension,
\eqref{eq:local-reconstruction-estimates},
\eqref{eq:supp-general-Garding}, and \eqref{eq:supp-general-core} provide the
same hypotheses.  The spectral-projection conclusion is the usual consequence
of self-adjoint norm-resolvent convergence for isolated spectral clusters.
\end{proof}

\subsection{Independent all-dimensional metric head}

For each chart choose a mesoscopic radius $\rho_h\downarrow0$ with
$h/\rho_h\to0$.  Connect nearby vertices by endpoint blocks whose edge length is the
length of the corresponding short physical curve in $g$.  Let $D_h^{\rm met}$ be
the direct sum of these source--terminal blocks.

\begin{theorem}[Exact finite metric head and geodesic convergence]
\label{thm:supp-global-metric}
For $f\in C(V_h)$,
\begin{equation}
{
L_h^{\rm met}(f)
=\|[D_h^{\rm met},\pi_h(f)]\|
=\max_{\{x,y\}\in E_h}\frac{|f(x)-f(y)|}{\ell_h(x,y)}.}
\label{eq:supp-edge-lip}
\end{equation}
The pure-state Connes distance is exactly the graph shortest-path metric.  On every
closed Riemannian manifold,
\begin{equation}
{
\sup_{x,y\in V_h}
|d_h^{\rm met}(x,y)-d_g(x,y)|\longrightarrow0.}
\label{eq:supp-global-metric-convergence}
\end{equation}
For smooth $f$, $L_h^{\rm met}(f|_{V_h})\to\|\nabla_gf\|_\infty$.
\end{theorem}

\begin{proof}
The commutator of one endpoint block has off-diagonal entry
$(f(y)-f(x))/\ell_h(x,y)$, so the operator norm gives
\eqref{eq:supp-edge-lip}.  Telescoping along graph paths gives one inequality in the
Connes dual formula; the distance-to-a-fixed-vertex function has seminorm one and
gives equality.

Every graph edge has the length of an actual curve, so
$d_g\le d_h^{\rm met}$.  Conversely, partition a minimizing geodesic into pieces of
length comparable with $\rho_h$ and replace the partition points by mesh vertices
within $O(h)$.  Bounded geometry gives
\[
d_h^{\rm met}(x,y)
\le d_g(x,y)+C\left(\rho_h+\frac h{\rho_h}\right).
\]
Choosing, for example, $\rho_h=h^{1/2}$ proves uniform distance convergence.  Taylor
expansion along the short edges gives the smooth seminorm convergence.
\end{proof}

For these endpoint metric triples the finite Lip-norm is exactly the
ordinary Lipschitz seminorm of the graph shortest-path distance.  The metric
on $S(C(V_h))$, identified with probability measures on $V_h$, is therefore
the Kantorovich transport distance for $d_h^{\rm met}$
\cite{Rieffel1999}.  If $p,q$ are probability measures supported on $V_h$,
the same transport plans give
\[
 \left|\operatorname{mk}_{L_h^{\rm met}}(p,q)-W_{1,d_g}(p,q)\right|
 \le\sup_{x,y\in V_h}|d_h^{\rm met}(x,y)-d_g(x,y)|.
\]
Here $W_{1,d_g}$ denotes the transport distance between the same measures
viewed on $M$.  Unlike the general quadratic graph seminorm, this particular
endpoint construction thus yields an immediate comparison on all states,
not only on vertices.

\begin{corollary}[Conditional stable compact-spin approximation]
\label{cor:supp-compact-spin-image}
A closed marked Riemannian spin manifold admitting a compatible stable spin
discretization has a finite-dimensional approximation of its doubled spin
Dirac operator in generalized norm resolvent, together with convergent
independent metric heads.  Consequently, essential surjectivity of this
particular atlas-based spin construction follows if such compatible stable
discretizations are supplied for all target manifolds.
\end{corollary}

\begin{proof}
Apply Theorem~\ref{thm:supp-compact-spin-resolvent} to the spin operators and
Theorem~\ref{thm:supp-global-metric} to the independent metric heads.
The last assertion records explicitly the additional existence requirement.
\end{proof}

\section{Finite rational indefinite memory}
\label{app:rational-krein}
% =============================================================================

The positive Stieltjes theory in Subsection~\ref{subsec:dynamic-memory} is not the full dynamic
image.  A finite Krein/Pontryagin tail is instead controlled by the number of
negative squares of its generalized Nevanlinna kernel.

\begin{definition}[Normalized rational Hermitian dynamic function]
\label{def:supp-rational-Hermitian}
Let $H$ be finite dimensional and let $Q$ be an $\End(H)$-valued rational function
satisfying
\begin{equation}
Q(\bar z)^*=Q(z),
\qquad
\lim_{z\to\infty}Q(z)=0.
\label{eq:supp-rational-symmetry}
\end{equation}
Write its Laurent expansion at infinity as
\begin{equation}
Q(z)=-\sum_{n\ge0}M_nz^{-n-1},
\qquad M_n=M_n^*.
\label{eq:supp-rational-Laurent}
\end{equation}
Let $\mathcal K_Q$ be the span of Hankel columns
$c_j(h)=(M_{j+r}h)_{r\ge0}$ and define
\begin{equation}
[c_i(u),c_j(v)]_Q=\langle u,M_{i+j}v\rangle.
\label{eq:supp-Hankel-form}
\end{equation}
After quotienting the null space, let $A_Qc_j(h)=c_{j+1}(h)$ and
$\Gamma_Qh=c_0(h)$.
\end{definition}

\begin{theorem}[Complete normalized rational generalized-Nevanlinna realization]
\label{thm:supp-complete-rational-Krein}
The construction above is a minimal finite Pontryagin realization and
\begin{equation}
{Q(z)=\Gamma_Q^+(A_Q-z)^{-1}\Gamma_Q.}
\label{eq:supp-rational-realization}
\end{equation}
Its ordinary dimension is the stabilized coefficient-Hankel rank
\begin{equation}
\dim\mathcal K_Q=\McM(Q),
\label{eq:supp-rational-degree}
\end{equation}
and its negative index is the inertia of
\eqref{eq:supp-Hankel-form}.  The Nevanlinna kernel
\begin{equation}
N_Q(z,w)=\frac{Q(z)-Q(w)^*}{z-\bar w}
\label{eq:supp-Nevanlinna-kernel}
\end{equation}
has exactly that many negative squares.  Every other minimal finite
Pontryagin realization of $Q$ is related to this one by a unique source-fixing
Pontryagin unitary.
\end{theorem}

\begin{proof}
By construction, the output functional $L_Q$ on the Hankel column model obeys
$L_QA_Q^n\Gamma_Q=M_n$.  For large $z$,
\[
L_Q(A_Q-z)^{-1}\Gamma_Q
=-\sum_{n\ge0}\frac{M_n}{z^{n+1}}=Q(z),
\]
and rational continuation proves \eqref{eq:supp-rational-realization}.  The
power-cyclic span of $\Gamma_QH$ is all of $\mathcal K_Q$ by definition.  In finite
dimension, the resolvent-source span equals the power-cyclic span: expansion at
infinity gives one inclusion and a finite Vandermonde inversion at sufficiently many
regular spectral parameters gives the other.  Thus the realization is source
minimal.

The dimension is the stabilized Hankel rank by construction.  Choosing a Hilbert
majorant and a congruence which writes the nondegenerate Hankel form as a fundamental
symmetry $J_Q$ gives $A_Q^*J_Q=J_QA_Q$ and
$Q(z)=\Gamma_Q^*J_Q(A_Q-z)^{-1}\Gamma_Q$.  The standard resolvent identity gives
\[
N_Q(z,w)
=\Gamma_Q^*J_Q(A_Q-z)^{-1}(A_Q-\bar w)^{-1}\Gamma_Q.
\]
Every finite kernel Gram is therefore a compression of the Pontryagin form, so it
has at most its negative index.  Source minimality makes the resolvent-source vectors
span the whole state space and attains the full index.  Finally, if two
minimal realizations have the same $Q$, their complete Hankel moments agree.
Mapping $A_Q^j\Gamma_Qh$ to the corresponding cyclic vector in the second
realization preserves the indefinite form, is onto, fixes the source, and
intertwines the state operators.  Uniqueness follows from cyclicity.
\end{proof}

\begin{theorem}[Pole--Hankel invariants]
\label{thm:supp-pole-Hankel}
For a normalized rational Hermitian $Q$:
\begin{enumerate}[label=\textup{(P\arabic*)},leftmargin=3em]
\item the principal part at each pole determines the local McMillan degree and all
      Jordan block sizes through shifted principal-part Hankel ranks;
\item at each real pole, the principal-part Hankel Gram determines the complete
      local Pontryagin inertia;
\item every nonreal conjugate pole pair forms a neutral paired root space, and its
      upper-half-plane local McMillan degree contributes equally to positive and
      negative index;
\item the sum of the local contributions equals the minimum global negative index.
\end{enumerate}
\end{theorem}

\begin{proof}
At a pole $\lambda$, multiply by a sufficiently high power of $(z-\lambda)$ and
read the finite principal-part coefficient sequence.  The ranks of its shifted
Hankel matrices are the controllability/observability ranks of the nilpotent local
state operator and their successive differences recover the Jordan chain lengths.
For a real pole the local Hankel form is Hermitian and Sylvester inertia is invariant
under every minimal state-coordinate change, hence gives the local Pontryagin
signature.  Hermitian symmetry pairs a nonreal root space with its conjugate; their
indefinite form is neutral off-diagonal, so the two signs occur with equal
multiplicity.  The primary decomposition is Pontryagin orthogonal, and the global
index is the sum of local indices.
\end{proof}

A general finite rational Hermitian function has a unique split
\begin{equation}
Q(z)=P(z)+Q_0(z),
\qquad
\lim_{z\to\infty}Q_0(z)=0,
\label{eq:supp-rational-split}
\end{equation}
with a Hermitian polynomial $P$.

\begin{corollary}[All finite rational Hermitian dynamics]
\label{cor:supp-all-rational}
The strictly proper part $Q_0$ is realized by
\Cref{thm:supp-complete-rational-Krein}.  The nonconstant polynomial jet $P$ has a
canonical finite infinity-chain/regular descriptor realization.  Their primary
direct sum gives an equivalence, up to source-fixing descriptor Pontryagin unitary,
\begin{equation}
{
\begin{gathered}
\text{finite rational Hermitian dynamic functions}\\
\longleftrightarrow\\
\text{minimal regular Pontryagin descriptor realizations}.
\end{gathered}}
\label{eq:supp-all-rational-equivalence}
\end{equation}
This includes finite real and nonreal poles, arbitrary finite Jordan chains, and all
polynomial jets at infinity.
\end{corollary}

\begin{proof}
The pencil of any regular finite descriptor realization has a primary decomposition
at infinity and at finite spectral points.  On the finite primary block the transfer
function is strictly proper and is minimized by
\Cref{thm:supp-complete-rational-Krein}.  On the generalized zero eigenspace of an
anchored descriptor pencil, the resolvent expansion is polynomial and is represented
by the canonical nilpotent infinity chain.  The direct sum attains both minimum
state dimensions and indices.  Any second minimal realization splits into the
same invariant primary blocks, and the unique source-fixing unitaries on the finite
and infinity blocks combine to the unique descriptor equivalence.
\end{proof}

\begin{remark}[Positive OS memory is the index-zero branch]
\label{scope:supp-positive-index-zero}
A positive Stieltjes/OS memory packet has zero negative-square index.  A nonzero
physical Pontryagin index requires an independently occurring fundamental-symmetry
or signed packet.  Neither OS reflection, temporal orientation, charge conjugation,
nor a determinant/Pfaffian line determines that sign by itself.  The unrestricted
nonrational infinite-dimensional indefinite essential image is not claimed here.
\end{remark}


\section{Completed inverse fibres}\label{app:completed-fibres}
\subsection{Joint Gram realizations and the finite stratified fibre}

Let $E_1,\ldots,E_m$ be finite Hilbert spaces representing minimal marginal
realization spaces, normalized so that their marginal Gram operators are the
identity.  A joint realization is encoded by a block correlation operator
\begin{equation}
\mathcal K=[K_{ab}]_{a,b=1}^m,
\qquad
K_{aa}=I_{E_a},
\qquad
\mathcal K\succeq0.
\label{eq:supp-correlation-spectrahedron}
\end{equation}
Let $\Corr_X$ denote this correlation spectrahedron and let $\mathcal G_X$ be
the product of the unitary gauge groups of the marginal realization spaces.

\begin{theorem}[Joint Gram classification]
\label{thm:supp-common-source}
Every minimal joint realization of the retained marginal Gram blocks is
classified, up to a unitary on the common realization space, by a point of
$\Corr_X/\mathcal G_X$.  The support rank of $\mathcal K$ is the minimum
possible joint realization dimension.  A representative $\mathcal K$ has
residual gauge group exactly $\Stab_{\mathcal G_X}(\mathcal K)$.
\end{theorem}

\begin{proof}
A joint realization has a Gram factorization $\mathcal K=S^*S$, with the
identity diagonal expressing the chosen normalization of the marginal blocks.
Conversely every positive block operator with identity diagonal admits such a
factorization.  Restricting the realization space to $\overline{\Ran S}$ gives
the minimal realization, whose dimension is $\rank\mathcal K$ in the finite
case.  Two minimal Gram factorizations of the same positive operator are
related by a unique unitary on the common range.  Independent changes of
marginal frames act by block congruence, and the residual action is precisely
the stated stabilizer.
\end{proof}

\begin{proposition}[Pairwise positivity need not give a joint realization]
\label{cth:supp-pairwise-not-global}
Pairwise admissible Gram blocks need not assemble into one positive joint Gram
operator.  For example,
\begin{equation}
\begin{pmatrix}
1&-3/4&-3/4\\
-3/4&1&-3/4\\
-3/4&-3/4&1
\end{pmatrix}
\label{eq:supp-three-block-counterexample}
\end{equation}
has every $2\times2$ principal block positive but has eigenvalue $-1/2$.
Thus pairwise compatibility does not imply global Gram compatibility.
\end{proposition}

Let $\mathcal J_X$ be the current fibre, $\mathcal M_X$ the complete
dynamic-memory fibre over the fixed spectral data, $\mathcal P_X$ the retained
auxiliary-marking fibre, $\mathcal C_X^{\rm cal}$ the residual calibration
class, and $\mathcal R_X$ the marked monodromy representation space.  We call a
finite target \emph{fully reduced} when all null and spectator directions have
been quotiented out, the domain, weight, orientation, determinant-line, normalization, and auxiliary-marking
conventions are fixed, every retained Gram block is minimal, and a specified
finite family of bounded separating functionals (equivalently, finite-rank
defect operators) has trivial common kernel on any additional retained
direction.  This last condition is the standard separation hypothesis that
replaces the programme-specific ``zero-exhaustion'' terminology.

\begin{theorem}[Finite stratified inverse-fibre exhaustion]
\label{thm:supp-finite-exhaustion}
For a completed fully reduced finite target $\mathsf S^\sharp$,
\begin{equation}
{
\RReal_X^{\sharp,\min}(\mathsf S)
\cong
\bigsqcup_{[\mathcal K]\in\Corr_X/\mathcal G_X}
\frac{
\mathcal J_X\times\mathcal M_X\times\mathcal P_X
\times\mathcal C_X^{\rm cal}\times\mathcal R_X
}{\Stab_{\mathcal G_X}(\mathcal K)}.}
\label{eq:supp-finite-fibre}
\end{equation}
The factors vary, respectively, directed dynamics, dynamical memory,
auxiliary markings, calibration, and global bundle data while preserving
the fixed spectral packet.  In general the quotient is stratified and not a
Cartesian product.
\end{theorem}

\begin{proof}
The separating family in the definition of full reduction rules out any
additional independent retained direction after the null quotients have been
taken.  The graph inverse theorem classifies the directed-current freedom; the
positive and Pontryagin/descriptor theorems classify returned dynamic memory;
the retained auxiliary markings give $\mathcal P_X$; the calibration data give
$\mathcal C_X^{\rm cal}$; and marked bundle descent gives $\mathcal R_X$.
Theorem~\ref{thm:supp-common-source} classifies all compatible joint Gram
realizations and identifies their residual stabilizers.  Equivalence of two
fully reduced realizations is therefore simultaneous equivalence in these
factors under the common stabilizer, yielding \eqref{eq:supp-finite-fibre}.
Proposition~\ref{cth:supp-pairwise-not-global} explains why the general fibre
cannot be replaced by a product of independently quotiented factors.
\end{proof}

\subsection{Quasilocal and semifinite completion}

A completed finite-stage system consists of increasing finite-trace projections
$P_n\uparrow I$, fully reduced finite fibres $\mathfrak F_n$, bonding maps onto
their stable images, tightness of the finite realization and operator-measure data, compatible
bundle transport, compatible self-adjoint resolvent data, and a compatible
normal semifinite weight.  In addition, choose at each stage a family
$\mathscr F_n$ of bounded separating functionals (or finite-rank defect
operators), compatible under restriction, such that the cofinal family
$\bigcup_n\mathscr F_n$ separates every nonzero additional quasilocal retained
direction.  Equivalently, the common kernel of this cofinal separating family
on the quotient by the declared finite-stage coordinates is zero.

\begin{theorem}[Completed quasilocal/semifinite inverse-fibre theorem]
\label{thm:supp-semifinite-exhaustion}
For such a completed system, the retained inverse fibre is exactly
\begin{equation}
{
\RReal_\infty^{\sharp,\min}(\mathsf S)
\cong
\varprojlim_n
\left[
\bigsqcup_{[\mathcal K_n]\in\Corr_n/\mathcal G_n}
\frac{
\mathcal J_n\times\mathcal M_n\times\mathcal P_n
\times\mathcal C_n^{\rm cal}\times\mathcal R_n
}{\Stab(\mathcal K_n)}
\right].}
\label{eq:supp-semifinite-fibre}
\end{equation}
The inverse limit is taken over the stable images, is nonempty exactly when the
stable-image system is nonempty, and contains no additional retained coordinate
under the stated separation and tightness hypotheses.  On a factor-preserving
stratum it may be written as the product of the five component inverse limits
and the correlation inverse limit followed by the compatible gauge quotient;
without factor-preserving transport the stratified form
\eqref{eq:supp-semifinite-fibre} is essential.
\end{theorem}

\begin{proof}
Surjective restriction to stable images makes the inverse system of nonempty
compact finite strata satisfy the finite-intersection property, so compatible
threads exist whenever every stable image is nonempty.  Tightness prevents a
nonzero realization direction or dynamic mass from escaping every finite
projection.  The compatible resolvents determine one closed self-adjoint
operator through the resolvent identity, adjoint symmetry, injectivity, and
dense range.  The compatible finite-trace weights increase to a unique
lower-semicontinuous normal semifinite weight by monotone convergence, while
compatible bundle transport fixes the remaining global data.  Finally, if an additional
retained quasilocal coordinate existed, the cofinal separating family would
contain a functional or defect operator nonzero on it, contradicting the
trivial-common-kernel hypothesis.  The finite fibres are given by
Theorem~\ref{thm:supp-finite-exhaustion}; taking compatible threads proves
\eqref{eq:supp-semifinite-fibre}.
\end{proof}


\end{document}